% ===================================================================
% FST-I: Thermodynamische Stabilit\"at fundamentaler Parameter
% Zieljournal: Foundations of Physics
% Version: 3.0 (Quellencheck-Revision)
% ===================================================================

\documentclass[12pt,a4paper]{article}

% Packages
\usepackage[utf8]{inputenc}
\usepackage[ngerman]{babel}
\usepackage[T1]{fontenc}
\usepackage{lmodern}
\IfFileExists{glyphtounicode.tex}{\input{glyphtounicode}}{}
\pdfgentounicode=1
\usepackage{amsmath,amssymb,amsthm}
\usepackage{physics}
\usepackage{graphicx}
\usepackage{booktabs}
\usepackage[margin=2.5cm]{geometry}
\usepackage{natbib}
\usepackage{enumitem}
\usepackage[expansion=false]{microtype}
\usepackage{tabularx}
\usepackage{array}
\usepackage[unicode=true,pdfencoding=auto,psdextra]{hyperref}
\usepackage{cleveref}
\renewcommand*{\HyperDestNameFilter}[1]{de.#1}

% Theorem environments
\newtheorem{proposition}{Proposition}
\newtheorem{conjecture}{Conjecture}
\newtheorem{definition}{Definition}
\newtheorem{assumption}{Annahme}
\newtheorem{lemma}{Lemma}
\theoremstyle{remark}
\newtheorem{remark}{Bemerkung}
\newcolumntype{Y}{>{\raggedright\arraybackslash}X}
\setlength{\emergencystretch}{2em}

% Title
\title{Funktionale Stabilitätstheorie I: Thermodynamische Stabilität fundamentaler Parameter}

\author{Lukas Geiger\\
\textit{Unabh\"angiger Forscher, Bernau, Deutschland}\\
ORCID: \url{https://orcid.org/0009-0005-7296-1534}}

\date{M\"arz 2026}

\hypersetup{
  pdftitle={Funktionale Stabilitätstheorie I: Thermodynamische Stabilität fundamentaler Parameter},
  pdfauthor={Lukas Geiger},
  pdfsubject={Thermodynamische Stabilität fundamentaler Parameter in einem spieltheoretischen Rahmen},
  pdfkeywords={Spieltheorie, Entropieproduktion, Standardmodell, Fine-Tuning, Nash-Gleichgewicht, Quantenmechanik, Symmetriebrechung, Envariance, Dekohärenz}
}

\begin{document}

\maketitle

\begin{abstract}
Dieser Beitrag schl\"agt einen spieltheoretischen interpretativen Rahmen f\"ur das Standardmodell der Teilchenphysik vor, der auf thermodynamischen Optimierungsprinzipien basiert. Wir untersuchen die Hypothese, dass der beobachtete Teilcheninhalt und die Kopplungskonstanten konsistent mit einem Nash-Gleichgewicht eines mehrskaligen thermodynamischen Spiels sind, wobei die Payoff-Funktion die integrierte Entropieproduktion \"uber kosmologische Zeitskalen ist. Innerhalb dieses Rahmens wird die Quantenmechanik als Gleichgewicht gemischter Strategien uminterpretiert, wobei die Born-Regel \"uber umgebungsgest\"utzte Invarianz (Envariance) als m\"ogliche Gleichgewichtsverteilung umgedeutet wird---eine konzeptuelle Uminterpretation, keine neue Herleitung. Symmetriebrechung wird als \"Ubergang zwischen Gleichgewichten gelesen, und Fine-Tuning wird als m\"ogliche eingeschr\"ankte Optimierung statt als Zufall uminterpretiert. Im Unterschied zu bisherigen thermodynamischen Ans\"atzen zum Fine-Tuning liefert FST-I einen hierarchischen mehrskaligen Rahmen, in dem Parameterselektions-Hypothesen \"uber Stabilit\"at in gekoppelten Sektoren gepr\"uft werden. Der Ansatz modifiziert nicht den mathematischen Formalismus der Quantenfeldtheorie; er liefert keine quantitativen Vorhersagen jenseits der etablierten Physik, bietet aber ein vereinigendes interpretatives Vokabular, das in stochastischer Spieltheorie und Mean-Field-Dynamik begr\"undet ist. Ein semi-analytisches Sternmodell zeigt, dass die beobachtete Feinstrukturkonstante innerhalb eines breiten Plateaus nahezu maximaler integrierter stellarer Entropieproduktion liegt ($\mathcal{S}/\mathcal{S}_{\max} = 0{,}988$), was eine erste quantitative Konsistenzpr\"ufung---wenn auch noch keinen entscheidenden Test---f\"ur die Entropie-Optimalit\"atshypothese liefert. Das Plateau selbst (mit $\mathcal{S}/\mathcal{S}_{\max} > 0{,}92$ \"uber den gesamten lebensfähigen Bereich) zeigt, dass stellare Entropieproduktion allein $\alpha$ nicht scharf selektiert; ein zweidimensionaler Scan \"uber $\alpha$ und $m_e/m_p$ ergibt $\mathcal{S}/\mathcal{S}_{\max}^{2D} = 0{,}47$, was zeigt, dass eine uneingeschr\"ankte einzelskalige stellare MEPP kein eigenst\"andiges Selektionsprinzip ist und dass hierarchische mehrskalige Einschr\"ankungen aus Atom- und chemischer Physik eine unabh\"angig pr\"ufbare Arbeitshypothese bleiben. Wir formulieren Falsifizierbarkeitskriterien und identifizieren die Bedingungen f\"ur einen vollst\"andigen Mehrparameter-Test. Der Rahmen befindet sich im Stadium einer strukturierten Hypothese, die rechnerische Validierung erfordert, nicht einer pr\"adiktiven Theorie.
\end{abstract}

\textbf{Schl\"usselw\"orter:} Spieltheorie, Entropieproduktion, Standardmodell, Fine-Tuning, Nash-Gleichgewicht, Quantenmechanik, Symmetriebrechung, Envariance, Dekoh\"arenz

\clearpage

% ===================================================================
\section{Einleitung: Das Fine-Tuning-Problem der fundamentalen Physik}
\label{sec:introduction}
% ===================================================================

Das Streben nach einem einheitlichen theoretischen Rahmen in der fundamentalen Physik hat sich traditionell auf axiomatische Ans\"atze gest\"utzt, bei denen die Kernprinzipien der Quantenmechanik und die Parameterwerte des Standardmodells als irreduzible Gegebenheiten akzeptiert werden. Die historische Trennung zwischen dem mikroskopischen Quantenbereich und dem makroskopischen thermodynamischen Verhalten hat zu tiefgreifenden Interpretationsherausforderungen gef\"uhrt, insbesondere dem Messproblem und dem unerkl\"arlichen Fine-Tuning kosmologischer Konstanten \citep{Adams2019}.

Das Standardmodell enth\"alt 19 freie Parameter\footnote{Drei Eichkopplungen ($\alpha_{EM}$, $\alpha_s$, $\alpha_W$), sechs Quark-Massen, drei geladene Lepton-Massen, vier CKM-Parameter (drei Mischungswinkel und eine CP-verletzende Phase), die Higgs-Masse, das Higgs-VEV und der QCD-Vakuumwinkel $\theta_{\text{QCD}}$. Unter Einbeziehung von Neutrinomassen und PMNS-Mischung steigt die Anzahl auf $\sim$26.}---Kopplungskonstanten, Massen und Mischungswinkel---deren Werte experimentell bestimmt werden m\"ussen, anstatt aus ersten Prinzipien hergeleitet zu werden. Diese Werte zeigen bemerkenswertes ``Fine-Tuning'': kleine Variationen w\"urden Kernbindung verhindern, Atome destabilisieren oder die Massenhierarchie grundlegend ver\"andern. Diese Beobachtung hat eine umfangreiche Diskussion ausgel\"ost, wobei vorgeschlagene Erkl\"arungen vom anthropischen Prinzip \citep{Weinberg1987,Susskind2005} bis zu Landscape-Szenarien in der Stringtheorie \citep{Bousso2000} reichen.

Der vorgeschlagene Rahmen, die Funktionale Stabilit\"atstheorie~I (FST-I), bietet eine neue interpretative Perspektive. Er untersucht die Hypothese, dass die axiomatische Struktur der Quantenmechanik und das hochspezifische Fine-Tuning des Standardmodells als die emergenten, eingeschr\"ankten Gleichgewichtsergebnisse eines mehrskaligen Optimierungsspiels verstanden werden k\"onnen, das durch thermodynamische Prinzipien bestimmt wird. FST-I leitet keine Parameterwerte aus ersten Prinzipien ab und generiert keine quantitativen Vorhersagen jenseits der etablierten Physik; sein Beitrag ist ein vereinigendes konzeptuelles Vokabular, das bekannte Parametersensitivit\"aten innerhalb einer spieltheoretischen Architektur organisiert.

In diesem Paradigma wird das Universum als ein komplexes, nicht-gleichgewichtiges thermodynamisches System modelliert, dessen gro\ss{}skalige Dynamik konsistent mit hohen Entropieproduktionsraten \"uber kosmologische Epochen hinweg ist \citep{MEPP2004}. Die effektiven Freiheitsgrade in diesem Spiel sind die Vakuumfeldsektoren, deren Kopplungen die thermodynamischen Payoffs bestimmen, die durch lokale Entropieproduktion definiert sind. Durch die Verkn\"upfung von nicht-kooperativer Spieltheorie, umgebungsinduzierter Superselektion (Einselection) und dem Maximum Entropy Production Principle (MEPP) bietet FST-I einen vereinigenden interpretativen Rahmen f\"ur den Ursprung von Quantenwahrscheinlichkeiten, die Notwendigkeit spontaner Symmetriebrechung und die Lage der physikalischen Konstanten im Parameterraum \citep{QuantumOptimization,LasryLions2007,HiggsVEV2024}.

% ===================================================================
\section{Strukturelle Ausrichtung: FST-I als Pattern-A-Instanz}
\label{sec:structural-alignment}
% ===================================================================

Die spieltheoretische Rahmung des Standardmodells in dieser Arbeit ist eine vorgeschlagene Instanz der \textbf{Pattern-A-Stabilit\"atslogik} im FST-Forschungs\"okosystem (siehe \cite{FST-Unified2026}), kein Beweis, dass das Standardmodell die Pattern-A-Kriterien bereits erf\"ullt. Pattern~A bezeichnet eine wiederkehrende Dreischrittstruktur: (i)~ein unendlich-dimensionaler Konfigurationsraum wird durch Einschr\"ankungen auf eine endliche effektive Dimension reduziert, (ii)~verbleibende Instabilit\"aten signalisieren verborgene Freiheitsgrade, die Eichfixierung erfordern, und (iii)~das resultierende eingeschr\"ankte Funktional ist positiv-definit, was Stabilit\"at zertifiziert. Die j\"ungste \textbf{eichtheoretische Reformulierung der Riemannschen Vermutung} (siehe \cite{GeigerRH2026c}) liefert das strukturelle Vorbild f\"ur diese Architektur.

Innerhalb dieses verallgemeinerten Stabilit\"atsrahmens wird das Standardmodell-Gleichgewicht wie folgt verstanden:
\begin{enumerate}[label=\alph*)]
    \item \textbf{Analytische Rang-Endlichkeit (Null-Tail):} Der unendlich-dimensionale Hilbert-Raum der Feldkonfigurationen bleibt handhabbar, weil der ``strategisch gangbare'' Unterraum analytisch begrenzt ist. Der hochenergetische (UV-)Tail wird durch die Entropieproduktions-Einschr\"ankung neutralisiert, was Informationslecks in unphysikalische Sektoren verhindert.
    \item \textbf{Endliche Negativit\"at als Eichsignal:} Die Existenz ``ungebrochener'' oder masseloser Symmetrien ist das Analogon der endlichen Negativit\"at im arithmetischen Kern. Diese sind keine Defekte, sondern Signale verborgener Freiheitsgrade, die durch eine Eichverschiebung stabilisiert werden m\"ussen---dargestellt hier durch Symmetriebrechung und den Higgs-Mechanismus.
    \item \textbf{Eingeschr\"ankte funktionale Positivit\"at (Nash-Stabilit\"at):} In der vorgeschlagenen Lesart w\"are das Nash-Gleichgewicht (Standardmodell) der ``eichstabile'' Zustand, in dem das Funktional der freien Energie ein eingeschr\"anktes Minimum erreicht. So wie die RH auf eine Rang-2-Stabilisierung reduziert wird, m\"usste Standardmodell-Stabilit\"at durch eine endliche Menge von Parameterkopplungen verifiziert werden, die globale thermodynamische Positivit\"at gew\"ahrleisten.
\end{enumerate}

Diese Ausrichtung transformiert FST-I von einer plausiblen Analogie in eine \emph{strukturelle} Analogie des ``Funktionale Positivit\"at unter Einschr\"ankung''-Mechanismus. Wir betonen, dass dies eine Analogie bleibt---der formale Beweis, dass der Standardmodell-Parameterraum einen resolventenstabilen Fixpunkt im Sinne von \citep{GeigerRH2026b} zul\"asst, wurde nicht erbracht. Der Wert der Analogie liegt in der Organisation der bekannten Parametersensitivit\"aten (Abschnitte~\ref{sec:finetuning}--\ref{sec:predictions}) innerhalb einer einheitlichen konzeptuellen Architektur. Eine erste quantitative Konsistenzpr\"ufung liefert das semi-analytische Sternmodell in Abschnitt~\ref{sec:stellar_model}, wo die beobachtete Feinstrukturkonstante 98,8\% der maximalen stellaren Entropieproduktion erreicht (Tabelle~\ref{tab:alpha_scan}).

\paragraph{Terminologie.} Im gesamten Beitrag bezeichnet \emph{resolvent-stabil}, in struktureller Analogie zur spektralen L\"uckeneigenschaft in \citep{GeigerRH2026b}, einen Fixpunkt der Dynamik, bei dem die St\"orungsanalyse zweiter Ordnung (Hessian oder Resolvent-Entwicklung) in allen Nicht-Eich-Richtungen strikt positive Eigenwerte liefert. Diese Terminologie wird als konzeptuelle Kurzbezeichnung verwendet; der formale Transfer auf physikalische Systeme erfordert in jedem Fall dom\"anenspezifische Verifikation, die f\"ur den Standardmodell-Parameterraum noch nicht durchgef\"uhrt wurde.

% ===================================================================
\section[Grundlagen der stochastischen Spieltheorie in physikalischen Systemen]{Grundlagen der stochastischen Spieltheorie\\in physikalischen Systemen}
\label{sec:foundations}
% ===================================================================

Um die Anwendung der Spieltheorie auf fundamentale Felder zu verstehen, muss man von klassischen Modellen mit endlich vielen Agenten zu kontinuierlichen, stochastischen Mean-Field-Spielen \"ubergehen.

\subsection{Vom Nash-Gleichgewicht zu Mean-Field-Spielen}

In der klassischen Mikro\"okonomie und Entscheidungstheorie beschreibt ein Nash-Gleichgewicht \citep{Nash1950} einen Zustand, in dem kein Agent einseitig von seiner gew\"ahlten Strategie abweichen kann, um seinen Payoff zu verbessern, gegeben die Strategien aller anderen Agenten. Bei der Erweiterung auf Quantensysteme erweist sich der klassische spieltheoretische Formalismus als ein eingeschr\"ankter Spezialfall einer umfassenderen Quanten-Spieltheorie \citep{QuantumOptimization}, in der die von-Neumann-Entropie und die von-Neumann-Morgenstern-Nutzenfunktion eine gemeinsame mathematische Struktur aufweisen (beide sind reellwertige, konkave Funktionale auf konvexen Zustandsr\"aumen) und der Hamilton-Operator als Kodierung der strategischen Payoff-Struktur interpretiert werden kann \citep{Abel2023,QuantumOptimization}.

Wenn die Anzahl der interagierenden Agenten (Feldanregungen) gegen unendlich geht ($N \to \infty$), geht der diskrete Nash-Gleichgewichts-Rahmen in ein Mean Field Game (MFG)-Gleichgewicht \"uber \citep{LasryLions2007}. In diesem Limes optimieren individuelle Agenten ihre Kostenfunktionale gegen\"uber der empirischen Verteilung der gesamten Population, getrieben durch kontrollierte McKean-Vlasov-artige stochastische Differentialgleichungen \citep{LasryLions2006}. Neuere Arbeiten haben verfeinerte $L^\infty$-Approximationsschranken f\"ur Nash-Gleichgewichte in Wasserstein-Raum-Formulierungen von MFGs etabliert \cite{DjeteTouzi2026}, was einen rigorosen metrischen Rahmen f\"ur die Konvergenz endlich-vieler-Agenten-Gleichgewichte zu ihren Mean-Field-Limiten liefert.

Das Ma\ss{} $m(p,t)$ wird als Wahrscheinlichkeitsverteilung \"uber Parameterkonfigurationen interpretiert, die die relative Stabilit\"at dynamischer Realisierungen kodiert---nicht als ontologisches Bekenntnis zu einem Multiversum. Der MFG-Rahmen bleibt rein mathematisch; physikalischer Gehalt ergibt sich erst durch die Spezifikation des Lagrangians~$L$ und der Wertfunktion~$V$.

\subsection{Hamilton-Jacobi-Bellman-Dynamik}

Die Evolution dieser Systeme beruht auf dynamischer Programmierung und Fixpunktargumenten, konkret gesteuert durch die Hamilton-Jacobi-Bellman-Gleichung:
\begin{equation}
\frac{\partial V}{\partial t} + \mathcal{H}\left(x, \nabla V, t\right) = 0
\end{equation}
wobei die Wertfunktion $V$ und der Kontroll-Hamiltonian $\mathcal{H}$ die optimalen R\"uckkopplungsstrategien bestimmen \citep{LasryLions2007}. Dieser stochastische Spielrahmen bildet das mathematische Fundament f\"ur die Interpretation von Quantenzust\"anden nicht als physikalische Wellen, sondern als Wahrscheinlichkeitsverteilungen \"uber einen komplexen Strategieraum.

\subsection{Quanten-Strategie-Korrespondenz}

\paragraph{Klassische Einbettung.}
Sei $\Delta_n = \{p \in \mathbb{R}_{\geq 0}^n : \sum_i p_i = 1\}$ der Simplex der klassischen gemischten Strategien \"uber $n$ reine Strategien und sei $\mathcal{D}(\mathcal{H})$ die Menge der Dichtematrizen auf $\mathcal{H} = \mathbb{C}^n$. Es gibt eine kanonische affine Einbettung $\iota: \Delta_n \hookrightarrow \mathcal{D}(\mathcal{H})$ gegeben durch $\iota(p) = \sum_i p_i |i\rangle\langle i|$, deren Bild genau die kommutative (diagonale) Teilmenge von $\mathcal{D}(\mathcal{H})$ ist. F\"ur jede diagonale Observable $A = \sum_i a_i |i\rangle\langle i|$ erf\"ullt der Erwartungswert $\mathrm{tr}(\iota(p) A) = \sum_i p_i a_i$, was exakt die klassische erwartete Payoff-Struktur reproduziert.

\paragraph{Nichtkommutative Erweiterung.}
Dichtematrizen stellen somit eine \emph{nichtkommutative Erweiterung} klassischer gemischter Strategien dar: diagonale Zust\"ande reproduzieren klassische Randomisierung, w\"ahrend Nichtdiagonal-Terme basisabh\"angige Interferenzstruktur kodieren, die keine klassische Simplex-Darstellung besitzt. Nichtdiagonal-Elemente $\rho_{ij}$ ($i \neq j$) repr\"asentieren phasensensitive Kopplungen zwischen Basisalternativen unter unit\"arer Evolution und nichtkommutierenden Observablen---sie sind \emph{keine} klassischen Korrelationen zwischen reinen Strategien \citep{Lindblad1976}.

\paragraph{Status und MFG-Kompatibilit\"at.}\mbox{}\\
Wir behandeln diese Quanten-Strategie-\allowbreak Korrespondenz als \emph{kommutative Spezialisierung}, nicht als Isomorphismus von Zustandsr\"aumen. Pr\"azise gilt: Unter einer Lindblad-Optimal-\allowbreak Transport-Geometrie auf~$\mathcal{D}(\mathcal{H})$, deren Generator die diagonale Unteralgebra invariant l\"asst und auf den Diagonalen einen Markov-Generator~$Q$ induziert, werden drei Eigenschaften erwartet: (i)~Diagonale Anfangszust\"ande bleiben unter der Vorw\"artsdynamik diagonal; (ii)~die nichtkommutative OT-Metrik reduziert sich auf den Diagonalen auf die diskrete Maas--Mielke--Erbar OT-Metrik auf~$\Delta_n$; (iii)~das gekoppelte Vorw\"arts-R\"uckw\"arts-System reduziert sich auf die endlichdimensionale Lasry--Lions Master-Gleichung. Diese Eigenschaften werden als Konjektur formuliert. Ein rigoroser Beweis unter den genannten geometrischen Bedingungen wurde nicht publiziert und stellt ein offenes mathematisches Problem dar (siehe auch Tabelle~\ref{tab:status-claims}). Falls etabliert, w\"aren klassische gemischte Strategien auf~$\Delta_n$ ein kommutativer Spezialfall eines Quanten-MFG auf~$\mathcal{D}(\mathcal{H})$.

Allerdings kann eine echte Quantenkorrespondenz mit Koh\"arenzen (Off-Diagonal-Elementen) \emph{nicht} aus klassischem~$W_2$-MFG allein gewonnen werden. Jede unit\"are Komponente $\partial_t \rho = -i[H, \rho] + \cdots$ mit nicht-diagonalem~$H$ treibt die Dynamik sofort von der Diagonaluntermannigfaltigkeit weg, und Geod\"atische in Quanten-OT erhalten~$\iota(\Delta_n)$ generisch nicht. Koh\"arenzen erfordern eine nichtkommutative Erweiterung des Zustandsraums und der Analysis, die \"uber die klassische Lasry--Lions-Theorie hinausgeht.

Der Rest dieses Beitrags baut auf der kommutativen Korrespondenz als konzeptuellem Organisationsprinzip auf. Die Erweiterung auf Koh\"arenzen bleibt ein offenes Problem f\"ur zuk\"unftige Arbeit.

% ===================================================================
\section{Quantenmechanik als gemischte Strategie}
\label{sec:quantum}
% ===================================================================

Um das thermodynamische Spiel mathematisch zu begr\"unden, bilden wir die formalen Komponenten der Quantenmechanik auf die Elemente der Spieltheorie ab. Im vorgeschlagenen Rahmen wird die axiomatische Struktur der Quantenmechanik uminterpretiert---auf der Ebene des Vokabulars, nicht neuer Physik---als thermodynamisches Optimierungsspiel, dessen Freiheitsgrade die Vakuumfeldsektoren sind. Diese Abbildung ist ein \emph{konzeptuelles Organisationsprinzip}: sie modifiziert nicht den Formalismus, leitet keine neuen Ergebnisse ab und generiert keine Vorhersagen jenseits der Standard-Quantentheorie (siehe Abschnitt~\ref{sec:foundations}, ``Geltungsbereich und formale Abgrenzung''). Der Wert liegt in der Aufdeckung struktureller Parallelen \"uber Skalen hinweg, nicht in der Ersetzung etablierter Formalismen.

\subsection[Zust\"ande als Strategien, Superposition als gemischte Strategie]{Zust\"ande als Strategien,\\Superposition als gemischte Strategie}

Wir interpretieren Quantenzust\"ande im Vokabular dieses thermodynamischen Spiels als Strategien. Ein reiner Quantenzustand $|\psi_i\rangle$ entspricht einer reinen Strategie. Folglich wird eine Quantensuperposition, definiert als $\sum_i c_i |\psi_i\rangle$, als Analogon einer gemischten Strategie behandelt. Das Betragsquadrat der Amplitude, $|c_i|^2$, repr\"asentiert das Wahrscheinlichkeitsgewicht, das dieser spezifischen Strategie zugewiesen wird. Diese Identifikation operiert innerhalb der in Abschnitt~\ref{sec:foundations} etablierten kommutativen Spezialisierung: die diagonale Dichtematrix $\rho = \sum_i |c_i|^2 |i\rangle\langle i|$ reproduziert die klassische gemischte-Strategie-Struktur, w\"ahrend Nichtdiagonal-Koh\"arenzen eine nichtkommutative Erweiterung erfordern, die \"uber den Umfang der vorliegenden Analyse hinausgeht.

Im mathematischen Formalismus stochastischer Spieldynamik erlaubt eine gemischte Strategie einem Agenten, Unsicherheit zu navigieren, indem er eine Wahrscheinlichkeitsverteilung \"uber den verf\"ugbaren Aktionsraum aufrechterh\"alt \citep{BusemeyerBruza2012}. In der hier entwickelten Analogie entspricht Superposition der gleichzeitigen Verf\"ugbarkeit mehrerer dynamischer Pfade, bevor Umgebungsinteraktion einen stabilen Record selektiert.

In dieser Lesart wird Superposition nicht als epistemische Unsicherheit \"uber eine verborgene zugrundeliegende Realit\"at behandelt, sondern als formale Koexistenz mehrerer Feldkonfigurationen vor der Payoff-Realisierung.

\subsection{Die Born-Regel als Gleichgewichtsverteilung}

Innerhalb dieses spieltheoretischen Konstrukts wird die Born-Regel nicht als Postulat der Standard-Quantentheorie ersetzt. Stattdessen kann sie als Gleichgewichtsverteilung der Strategien in einer Nash-Gleichgewichts-Analogie gemischter Strategien uminterpretiert werden (siehe Abschnitt~\ref{sec:quantum} f\"ur den pr\"azisen Umfang dieser Aussage). Der Payoff f\"ur das System wird als lokale Entropieproduktion zum Zeitpunkt der Messung (Interaktion) modelliert. Strategien mit einem h\"oheren erwarteten thermodynamischen Payoff w\"urden dann mit gr\"o\ss{}erer H\"aufigkeit realisiert. In dieser Lesart ist die durch $|\psi|^2$ definierte Wahrscheinlichkeitsverteilung die Kandidatin f\"ur die station\"are Verteilung des Spiels.

\subsubsection{Neuinterpretation via Envariance}

\textbf{Klarstellung:} Das Folgende stellt keine unabh\"angige Herleitung der Born-Regel dar. Es ist eine \emph{Neuinterpretation} von Zureks Envariance-Programm \citep{Zurek2003,Zurek2005} im spieltheoretischen Vokabular von FST-I. Zureks originale Herleitung ist selbst umstritten: Schlosshauer \& Fine (2005) \citep{SchlossFine2005} erheben Einw\"ande gegen die Behauptung, dass die Born-Regel allein aus Envariance ohne Zirkularit\"at folgt. Der Beitrag von FST-I hier besteht darin, zu beobachten, dass Zureks Symmetrieargumente nat\"urlich auf das Nash-Gleichgewichts-Konzept abbilden---eine konzeptuelle Verbindung, kein neues mathematisches Ergebnis.

Die Gleichgewichtsverteilungs-Lesart wird durch den Mechanismus der umgebungsgest\"utzten Invarianz, oder ``Envariance'', illustriert, die von Zurek entwickelt wurde \citep{Zurek2003,Zurek2005}.

Wenn ein Quantensystem $\mathcal{S}$ mit seiner Umgebung $\mathcal{E}$ verschr\"ankt wird, kann der zusammengesetzte reine Zustand \"uber die Schmidt-Zerlegung ausgedr\"uckt werden:
\begin{equation}
|\Psi_{SE}\rangle = \sum_k \alpha_k |s_k\rangle \otimes |e_k\rangle
\end{equation}
wobei $\{|s_k\rangle\}$ und $\{|e_k\rangle\}$ orthonormale Basen f\"ur die jeweiligen Hilbert-R\"aume $\mathcal{H}_S$ und $\mathcal{H}_E$ sind \citep{SchlossFine2005}.

Envariance besagt, dass eine lokale Eigenschaft des Systems unabh\"angig von Transformationen ist, die durch alleiniges Einwirken auf die Umgebung perfekt r\"uckg\"angig gemacht werden k\"onnen:

\begin{center}
\begin{tabularx}{\textwidth}{@{}>{\raggedright\arraybackslash}p{3.9cm}>{\raggedright\arraybackslash}p{3.8cm}Y@{}}
\toprule
\textbf{Transformation} & \textbf{Mechanismus} & \textbf{Implikation} \\
\midrule
Phasen-Envariance & $u_S = e^{i\beta_k}|s_k\rangle\langle s_k|$ & Phasen tragen keine strategische Information \\
Tausch-Envariance & Vertauschen von $|s_1\rangle$ und $|s_2\rangle$ & Gleiche Amplituden $\Rightarrow$ gleiche Wahrscheinlichkeiten \\
Z\"ahlungserweiterung & Feinstrukturierung der Umgebung & Verallgemeinerung auf alle rationalen $p_k$ \\
\bottomrule
\end{tabularx}
\end{center}

Zureks Envariance-Programm argumentiert unter seinen Symmetrie- und Feinstrukturierungsannahmen f\"ur $p_k \propto |\alpha_k|^2$. In FST-I wird dieser Weg als m\"ogliche Fixpunktsignatur in Analogie zu einem Nash-Gleichgewicht gelesen, nicht als neuer Beweis der Born-Regel. Da die Umgebung das System kontinuierlich \"uberwacht, werden Verteilungen au\ss{}erhalb der Born-Form als unter dekoh\"arenzartiger Relaxation instabil interpretiert, nicht als w\"ortliche Strategien mit gemessenen thermodynamischen Payoffs. Ein verwandter Relaxationsmechanismus---allerdings basierend auf fundamental anderer Dynamik---erscheint in der Pilotwellen-Interpretation: Valentini \& Westman \citep{ValentiniWestman2005} zeigen, dass Nicht-Born-Verteilungen durch Mischung durch die Pilotwelle in der de~Broglie--Bohm-Theorie zu $|\psi|^2$ relaxieren. Ihr Ergebnis st\"utzt die allgemeine Plausibilit\"at von Born-Regel-Relaxation \"uber Interpretationen hinweg, aber der zugrundeliegende Mechanismus (deterministische Hidden-Variable-Dynamik) unterscheidet sich vom hier vorgeschlagenen spieltheoretischen Bild.

\subsection{Das Pfadintegral als Summe \"uber Strategien}

Der konzeptuelle Ankerpunkt dieser Interpretation ist die Feynmansche Pfadintegral-Formulierung. Die \"Ubergangsamplitude wird durch die Summe \"uber alle m\"oglichen Historien gegeben:
\begin{equation}
A = \sum_{\text{paths}} e^{\frac{i}{\hbar} S[\text{path}]}
\end{equation}

Hier repr\"asentiert jeder einzelne Pfad eine vollst\"andige Strategie, und die Wirkung $S$ kodiert den kumulativen Payoff dieser Strategie \"uber die Zeit. Quanteninterferenz wirkt als Mechanismus des strategischen Wettbewerbs. Im semiklassischen Limes tragen nur Pfade nahe der station\"aren Phase konstruktiv bei.

\paragraph{Station\"are Phase als Sattelpunkt-Gleichgewicht.}
Eine wichtige Feinheit muss angesprochen werden: In der Minkowski-Signatur ergibt die station\"are-Phase-Bedingung $\delta S = 0$ einen \emph{Sattelpunkt}, kein Maximum. Das korrekte spieltheoretische Analogon ist daher nicht die Payoff-Maximierung eines einzelnen Agenten, sondern ein \emph{Nullsummen-Sattelpunkt-Gleichgewicht} (Minimax-Nash). In Hamilton-Form erster Ordnung,
\begin{equation}
S[q,p] = \int_{t_0}^{t_1} \left( p \cdot \dot{q} - H(q,p) \right) dt,
\end{equation}
sind die klassischen Gleichungen \"aquivalent zu den Sattelpunktbedingungen $\delta S / \delta q = 0$, $\delta S / \delta p = 0$---keine infinitesimale einseitige Variation in einer der konjugierten Variablen kann die entsprechende Seite des Sattelpunkts verbessern.

\paragraph{Euklidische Formulierung.}
Nach Wick-Rotation ($t \mapsto -i\tau$) wird das Gewicht zu $\exp(-S_E[q]/\hbar)$, und der semiklassische Limes wird durch \emph{Minima} von $S_E$ bestimmt. Eine w\"ortliche Payoff-Interpretation erh\"alt man dann durch Definition des Nutzens als $U[q] := -S_E[q]$, sodass Gleichgewichtspfade $U$ maximieren (\"aquivalent $S_E$ minimieren). Die Minkowski-station\"are-Phase-Aussage wird durch analytische Fortsetzung wiederhergestellt. Im Rest dieses Beitrags sollten Verweise auf ``Payoff-Maximierung'' in diesem euklidischen Sinne verstanden werden, sofern nicht anders angegeben.

Diese Formulierung ersetzt effektiv den klassischen Begriff einer einzelnen, eindeutigen Trajektorie durch ein Funktionalintegral \"uber eine Unendlichkeit quantenmechanisch m\"oglicher Trajektorien \citep{PathIntegral}. Auf Feldanregungen angewandt, dient die euklidische Wirkung als Nutzenfunktion \citep{PathIntegralOptimization}. Der gesamte Quantenpropagationsprozess ist eine parallelisierte mathematische Berechnung, bei der jede denkbare Strategie gleichzeitig evaluiert wird.

\subsection{Destruktive Interferenz als Strategieeliminierung}

Nicht-station\"are Pfade---solche, die weit vom Sattelpunkt der Wirkung entfernt liegen---erfahren schnelle Phasenoszillation, was zu destruktiver Interferenz f\"uhrt. Im hier verwendeten spieltheoretischen Vokabular ist destruktive Interferenz analog zur evolution\"aren Eliminierung: Beitr\"age fern vom station\"aren Sattel l\"oschen sich im Amplitudenraum aus. Die Analogie ist nur strukturell; sie macht aus Pfadintegral-Interferenz keine Populationsdynamik.

Diese Dynamik ist lose analog zu Selektionsprozessen in der evolution\"aren Biologie, insofern als nicht-optimale Strategien eliminiert werden. Die Analogie sollte jedoch nicht \"uberstrapaziert werden: destruktive Interferenz operiert auf komplexen Amplituden in einem einzelnen System, nicht auf Populationen mit Vererbung, Variation und differentieller Fitness. In einer hochkompetitiven stochastischen Umgebung weisen Strategien, die vom Sattelpunkt-Gleichgewicht abweichen, stark divergierende komplexe Phasen auf. Bei Aufsummierung im Pfadintegral summieren sich diese hochoszillierenden Terme zu Null, was effektiv die Wahrscheinlichkeit eliminiert, dass das System Nicht-Gleichgewichtszust\"ande einnimmt.

\subsection{Messung als Payoff-Realisierung}

Das kontroverse Konzept des Wellenfunktionskollapses wird als Payoff-Realisierung umgedeutet. Messung entspricht der Payoff-Realisierung durch irreversible Kopplung an makroskopische, thermodynamische Freiheitsgrade. Es gibt keinen mystischen ``physikalischen Kollaps''; vielmehr markiert das Messereignis das Ende der strategischen Flexibilit\"at. Das System legt sich auf einen definitiven Zustand fest und extrahiert den thermodynamischen Payoff.

Die Dekoh\"arenztheorie liefert den exakten physikalischen Mechanismus f\"ur diese Payoff-Realisierung \citep{MeasurementClassical}. Wenn ein Quantensystem mit einer makroskopischen Umgebung interagiert, werden die Phasenbeziehungen der Superposition irreversibel delokalisiert, wobei Information in die unbeobachtbaren Freiheitsgrade des umgebenden Bads transferiert wird. Dieser Prozess, Einselection, erzwingt ein effektives Verbot des Gro\ss{}teils des Hilbert-Raums, eliminiert fragile ``Schr\"odinger-Katzen''-Zust\"ande und l\"asst nur einen diskreten Satz robuster, klassischer Zeigerzust\"ande \"ubrig \citep{Zurek2003}.

Dar\"uber hinaus wird die Heisenbergsche Unsch\"arferelation als Erfordernis strategischer Flexibilit\"at verstanden: konjugierte Variablen k\"onnen nicht gleichzeitig fixiert werden, da dies den Strategieraum \"ubereinschr\"anken w\"urde und das System daran hindern w\"urde, ein optimales Gleichgewicht gemischter Strategien zu erreichen \citep{HeisenbergCovariant}.

\subsection{G\"ultigkeitsbereich und formale Abgrenzung}

Es ist wichtig, explizit anzugeben, was diese Interpretation beinhaltet und was sie vermeidet. Dieser Rahmen modifiziert nicht den mathematischen Formalismus der Quantenmechanik. Er f\"uhrt keine verborgenen Variablen ein, noch ver\"andert er die vorhersagenden Ergebnisse der Standard-Quantentheorie. Stattdessen liefert er eine funktionale Interpretation, in der Teile der axiomatischen Struktur der Quantenmechanik durch Optimierungs- und Gleichgewichtsprinzipien thermodynamischer Systeme organisiert werden.

\paragraph{Wert der spieltheoretischen Reformulierung.}
Wir beanspruchen keine neue Herleitung der Born-Regel oder neue quantitative Vorhersagen \"uber die Standard-Quantentheorie hinaus. Der Beitrag ist eine \emph{strukturelle Vereinigung}: Innerhalb von Zureks Envariance-Programm k\"onnen Born-Gewichte als Stabilit\"ats-/Fixpunktbedingung unter zul\"assigen Transformationen gelesen werden, in direkter Analogie zu KL/Freie-Energie-Gleichgewichten in der statistischen Mechanik (Gibbs-Minimum, FST-II) und Replikator-Dynamik in der evolution\"aren Biologie (Nash/ESS, FST-III). Das spieltheoretische Vokabular macht die impliziten Annahmen des Envariance-Ansatzes als ``Interaktionsregeln'' explizit---erlaubte Symmetrien, Informationszugang und Stabilit\"at unter einseitigen Abweichungen---und verbessert dadurch die Axiom-Transparenz. In diesem Sinne ist der Beitrag konzeptuelle Vereinigung und Axiom-Buchf\"uhrung, kein neues Theorem.

Durch die Beziehung ausgew\"ahlter Postulate der Quantenmechanik zu den Formalismen der stochastischen Spieldynamik und Envariance liefert FST-I eine gemeinsame Sprache, die Quantenmessung, chemische Selektion und biologische Evolution als Instanzen einer gemeinsamen Fixpunkt-/Variationsschablone verbindet \citep{SchlossFine2005}. Der mathematische Apparat der Quantenmechanik wird in seiner Gesamtheit bewahrt \citep{QuantumOptimization}, aber sein ontologischer Status wird neu beschrieben: Er wird als optimaler Kontrollrahmen f\"ur thermodynamische Effizienz unter Bedingungen unvollst\"andiger Information gelesen.

% ===================================================================
\section{Das Maximum Entropy Production Principle als globale Nutzenfunktion}
\label{sec:mepp}
% ===================================================================

Um von der Mikrodynamik quantenmechanischer gemischter Strategien zur Makrodynamik kosmologischer Evolution \"uberzugehen, m\"ussen wir die globale Nutzenfunktion, die das thermodynamische Spiel antreibt, explizit definieren. In FST-I wird diese Funktion durch das Maximum Entropy Production Principle (MEPP) bereitgestellt \citep{MEPP2004}.

\subsection{Formulierung von MEPP}

MEPP besagt in seiner starken Form, dass komplexe, nicht-gleichgewichtige Systeme, die stetigem \"au\ss{}erem Antrieb unterliegen, sich in Richtung thermodynamischer Pfade mit hoher Entropieproduktionsrate ($\dot{S}$) selbst organisieren k\"onnen \citep{MEPPFoundations}. Urspr\"unglich im Kontext der planet\"aren Klimatologie von Paltridge \citep{Paltridge1975} formuliert und nachfolgend \"uber Fluiddynamik, Chemie und Biologie verallgemeinert \citep{MEPP2004}, beschreibt MEPP eine vorgeschlagene Organisationstendenz zur Dissipation von Energiegradienten, kein allgemein akzeptiertes Variationsgesetz.

\subsection{Status und Kontroverse von MEPP}

\textbf{Wichtig:} MEPP ist ein aktives Diskussionsthema in der Gemeinschaft der Nicht-Gleichgewichtsthermodynamik. Sein Status als fundamentales Gesetz, als n\"utzliche Heuristik oder als rahmenabh\"angiges Postulat ist umstritten:

\begin{itemize}
    \item \textbf{Grinstein \& Linsker (2007)} zeigten, dass MEPP nicht allgemein f\"ur diskrete stochastische Systeme gilt und nicht ohne zus\"atzliche Annahmen aus der Standard-Nicht-Gleichgewichtsthermodynamik abgeleitet werden kann \citep{GrinsteinLinsker2007}.
    \item \textbf{Polettini (2013)} versch\"arfte diesen Einwand f\"ur Ziegler-MaxEP jenseits der linearen Antwort: Stochastisch-thermodynamische Modelle k\"onnen die von Ziegler geforderten Reziprozit\"atsrelationen verletzen; in Netzwerkmodellen folgen Steady States daher nicht schon aus MaxEP~\citep{Polettini2013}.
    \item \textbf{Dewar (2003)} argumentierte f\"ur MEPP als Konsequenz der Maximum-Entropie-Inferenz (Jaynes' Maximum-Entropie-Formalismus \citep{Jaynes1957} angewandt auf Pfadwahrscheinlichkeiten), aber diese Herleitung wurde aus technischen Gr\"unden kritisiert \citep{Dewar2003}.
    \item \textbf{Martyushev \& Seleznev (2006)} liefern die umfassendste \"Ubersicht und kommen zu dem Schluss, dass MEPP substanzielle empirische Unterst\"utzung \"uber viele Systeme hinweg hat, aber einer universell g\"ultigen Herleitung entbehrt \citep{MartyushevSeleznev2006}.
    \item \textbf{Sawada, Daigaku \& Toma (2025)} best\"atigten in einem aktuellen Review, dass MEPP makroskopische evolution\"are Trends gut beschreibt, eine rigorose Herleitung aus First Principles fern vom Gleichgewicht jedoch unvollst\"andig bleibt~\citep{Sawada2025}.
\end{itemize}

In FST-I wird MEPP als Arbeitshypothese behandelt---ein plausibles Organisationsprinzip mit empirischer Unterst\"utzung in mehreren Dom\"anen---statt als bewiesenes Axiom. Die Kernaussage des Papers erfordert nicht, dass MEPP universell gilt; sie erfordert, dass der Standardmodell-Parametersatz \emph{konsistent mit} einem eingeschr\"ankten Entropieproduktions-Optimum ist. Ob dies das Ergebnis von MEPP ist, das die Parameter \"uber kosmologische Zeitskalen ``selektiert'', oder eine strukturelle Eigenschaft stabiler Physik, die Falsifizierbarkeitskriterien in Abschnitt~\ref{sec:predictions} sind daf\"ur ausgelegt, dies zu entscheiden. Martyushev (2010) \citep{Martyushev2010} identifizierte zwei Grundfragen, die jede MEPP-Anwendung beantworten muss: (i)~ob das System lokales thermodynamisches Gleichgewicht erf\"ullt, und (ii)~ob die Entropieproduktion bilinear in Fl\"ussen und Kr\"aften ist. F\"ur die hier betrachteten Systeme fern vom Gleichgewicht ist Bedingung~(ii) in der Regel nicht erf\"ullt, was MEPP auf einen heuristischen Stabilit\"atsfilter statt eines rigorosen Selektionsprinzips einschr\"ankt.

\paragraph{Zirkularit\"atsvorbehalt.} Wenn MEPP als Payoff-Funktion verwendet wird und Nash-Gleichgewichte als MEPP-maximierende Konfigurationen definiert werden, wird die Behauptung, dass beobachtete Systeme Nash-Gleichgewichte sind, \emph{weil} sie Entropieproduktion maximieren, definitorisch zirkul\"ar. Der nicht-zirkul\"are Gehalt des Rahmens liegt in der \emph{Stabilit\"ats}-Aussage: Die beobachteten Parameter sind nicht blo\ss{} entropieproduzierend, sondern in zweiter Ordnung stabil gegen St\"orungen. Diese Stabilit\"atseigenschaft ist unabh\"angig von der Payoff-Definition testbar (Abschnitt~\ref{sec:finetuning}).

\subsection{Operationale Modi}

In multi-stabilen stochastischen Systemen manifestiert sich MEPP in zwei verschiedenen operationalen Modi \citep{MartyushevSeleznev2006}:
\begin{enumerate}
    \item \textbf{Zustandsselektionsprinzip:} Bestimmt, welcher spezifische station\"are Zustand in einer hochgradig entarteten Landschaft eingenommen wird;
    \item \textbf{Gradientenantwortprinzip:} Diktiert, wie das System auf St\"orungen reagiert.
\end{enumerate}

Auf die fundamentale Physik als FST-I-Hypothese \"ubertragen, legt MEPP nahe, dass die Regeln des Spiels---die fundamentalen Naturkonstanten---durch ihre F\"ahigkeit zu nachhaltiger thermodynamischer Dissipation eingeschr\"ankt sein k\"onnten \citep{EntropyMaximum}. Die Feldanregungen werden dann als Agenten modelliert, deren Wechselwirkungsquerschnitte und Kopplungsst\"arken gegen diese thermodynamische Nutzenfunktion gepr\"uft werden. Im FST-I-Rahmen impliziert dies nicht, dass physikalische Konstanten sich \"uber die Zeit \"andern; vielmehr werden die Konstanten als fr\"uhkosmologische Parameter behandelt, deren Stabilit\"atsstatus durch mehrskalige Constraints gefiltert werden k\"onnte.

\begin{remark}[MEPP als Stabilit\"atsfilter, nicht als Maximierungsprinzip]
\label{rem:mepp-stability-filter}
Eine entscheidende Verfeinerung ergibt sich aus der Resolventenanalyse zweiter Ordnung, die in der eichtheoretischen Reformulierung der Riemannschen Vermutung entwickelt wurde \citep{GeigerRH2026b}. Der FST-Rahmen beschreibt \emph{Stabilit\"atsprinzipien}, keine Maximierungsprinzipien. Die Dynamik f\"uhrender Ordnung ist entartet: viele Konfigurationen der Standardmodell-Parameter ergeben vergleichbare Entropieproduktionsraten erster Ordnung. Der vorgeschlagene Selektionsmechanismus w\"urde in zweiter Ordnung \"uber resolventenartige Kopplung zwischen entarteten und komplement\"aren Unterr\"aumen wirken. MEPP wird hier daher als \emph{Stabilit\"atsfilter-Heuristik} verwendet: Es motiviert die Suche nach eingeschr\"ankten, resolventenstabilen Fixpunkten innerhalb der entarteten Mannigfaltigkeit, nicht die Behauptung globaler Maximierung. Die beobachteten Standardmodell-Parameter sind keine globalen Entropieproduktions-Maxima; die offene FST-I-Behauptung ist, dass sie zweitordentlich stabilisierte Gleichgewichte sein k\"onnten.
\end{remark}

\paragraph{Skalentrennung und die FEP--MEPP-Dualit\"at.}
Ein verwandtes Problem ist die scheinbare Spannung zwischen Entropie-\emph{Minimierung} (Fristons Free Energy Principle, FEP) und Entropie-\emph{Maximierung} (MEPP). Im FST-Rahmen l\"ost sich diese Spannung durch Skalentrennung auf: Jeder Parametersektor minimiert lokal seine interne variationelle freie Energie und erh\"alt dadurch seine strukturelle Integrit\"at gegen Fluktuationen---das feldtheoretische Analogon einer Markov-Blanket-Struktur \citep{Friston2010}. Diese einzeln stabilen Sektoren bilden jedoch gemeinsam eine hocheffiziente dissipative Architektur, deren \emph{aggregierter} Effekt maximale Entropieproduktion auf der makroskopischen Skala ist. Interne FEP-Minimierung und externe MEPP-Maximierung sind daher keine konkurrierenden Prinzipien, sondern komplement\"are Aspekte desselben Mehrskalengleichgewichts; die formale Behandlung erfolgt in FST-III (Proposition~\texttt{fep-mepp}).

\begin{conjecture}[FEP--England-Dualit\"at]\label{lem:fep-england}
Sei $(I, B, E)$ eine Markov-Blanket-Zerlegung eines Systems im nichtgleichgewichtigen station\"aren Zustand, das lokale Detailed Balance erf\"ullt. Sei $\mathcal{A}_E$ die Menge zul\"assiger Blanket-Zustandstrajektorien \"uber $[0,T]$, und sei $\mathcal{Q}_I$ die Menge variationeller Dichten \"uber interne Zust\"ande. Dann gilt im Erwartungswert \"uber Trajektorien:
\[
  \max_{a \in \mathcal{A}_E} \langle \Delta S_{\mathrm{ext}} \rangle \;\Longleftrightarrow\; \min_{q \in \mathcal{Q}_I} F_{\mathrm{int}},
\]
wobei $\Delta S_{\mathrm{ext}} \geq 0$ die in den externen Zust\"anden produzierte Entropie ist, nach unten beschr\"ankt durch das Crooks-Fluktuationstheorem \citep{Crooks1999}, und $F_{\mathrm{int}} = \mathbb{E}_q[\ln q(s) - \ln p(s,o)]$ die variationelle freie Energie der internen Zust\"ande ist \citep{FristonFreeEnergy}. Die beiden Optimierungsprobleme operieren auf \emph{verschiedenen} Variablenmengen ($\mathcal{A}_E$ vs.\ $\mathcal{Q}_I$); die behauptete \"Aquivalenz besagt, dass am station\"aren Fixpunkt die L\"osung des einen Problems die L\"osung des anderen impliziert.
\end{conjecture}

Diese Dualit\"at wird ohne Beweis formuliert; das Argument ist eine Heuristik, motiviert durch das Crooks-Fluktuationstheorem, Englands dissipative Selbstreplikationsschranke \citep{England2013} und die variationelle Formulierung des Free Energy Principle. Eine rigorose Herleitung unter den genannten Bedingungen wurde nicht publiziert, und der genaue Sinn des Bikonditionals---ob es f\"ur alle station\"aren Zust\"ande oder nur f\"ur eine eingeschr\"ankte Klasse gilt---bleibt offen. Englands urspr\"ungliches Resultat liefert eine \emph{untere Schranke} f\"ur Entropieproduktion bei Selbstreplikatoren, keine \"Aquivalenz mit Freie-Energie-Minimierung; die hier vermutete Dualit\"at ist daher st\"arker als das, was die zitierte Literatur tr\"agt.

\begin{remark}
Diese Dualit\"at gilt im Erwartungswert \"uber Trajektorien, nicht pfadweise. Sie erfordert lokale Detailed Balance an der Blanket-Grenze und erstreckt sich nicht auf Systeme mit gebrochener Zeitumkehrsymmetrie. Die Dualit\"at ist daher als statistisches Ordnungsprinzip zu verstehen, nicht als dynamisches Gesetz.
\end{remark}

% ===================================================================
\section{Fine-Tuning als Optimierung}
\label{sec:finetuning}
% ===================================================================

Die These, dass das Standardmodell mit Entropieproduktions-Stabilit\"at kompatibel ist, muss operationalisiert werden, um testbare physikalische Konsequenzen zu liefern. Historisch wurden die exakten Werte der fundamentalen Konstanten als anthropische Koinzidenz betrachtet, fein abgestimmt auf einen willk\"urlichen Grad, um die Existenz von Beobachtern zu erm\"oglichen \citep{Adams2019}. FST-I testet eine aktivere Lesart: Die Parameterlandschaft k\"onnte um eingeschr\"ankten thermodynamischen Durchsatz organisiert sein, wobei biologische Komplexit\"at als lokalisiertes, hochdissipatives Nebenprodukt statt als prim\"ares Selektionsziel erscheint.

\subsection{Definition des Optimierungsfunktionals}

Um das Maximum Entropy Production Principle analytisch zu formulieren, definieren wir das totale Entropieproduktions-Funktional \"uber die relevante kosmologische Epoche:
\begin{equation}
\mathcal{S}[\vec{p}] = \int_{t_{BB}}^{t_\Lambda} \dot{S}(\vec{p}, t) \, dt
\end{equation}

Hier bezeichnet $\vec{p}$ den Satz dimensionsloser Standardmodell-Parameter, wie $\vec{p} = \{\alpha_{EM}, \alpha_s, \alpha_W, m_e/m_p, m_u/m_d, \ldots\}$. Die Integrationsgrenzen erstrecken sich vom Urknall ($t_{BB}$) bis zum Beginn der Dunkle-Energie-Dominanz ($t_\Lambda$), was die prim\"are Epoche der Strukturbildung und stellaren Nukleosynthese definiert.

Entscheidend ist, dass $\dot{S}$ nicht die grobk\"ornige Gesamtentropie des Universums darstellt (die von der Bekenstein-Hawking-Entropie supermassiver Schwarzer L\"ocher dominiert wird, die $\sim 10^{104}\,k_B$ aufweisen \citep{EntropyMaximum}). Schwarze L\"ocher sind thermodynamische Endpunkte, keine aktiven Produktionsmechanismen. Stattdessen ist $\dot{S}$ definiert als die baryonengewichtete Entropieproduktionsrate, die durch anhaltende nukleare, elektromagnetische und gravitationelle Prozesse erzeugt wird---prim\"ar angetrieben durch Sternfusion und nachfolgende Photonen-Thermalisierung.

\subsection{Variationsprinzip mit Einschr\"ankungen}

Wenn das Standardmodell das Ergebnis eines thermodynamischen Spiels ist, muss seine Parameterkonfiguration ein Variationsprinzip $\delta \mathcal{S} = 0$ mit $\delta^2 \mathcal{S} \leq 0$ (eingeschr\"anktes lokales Maximum) erf\"ullen, unter spezifischen Randbedingungen:
\begin{enumerate}
    \item \textbf{Nukleare Stabilit\"at:} Die Existenz gebundener, stabiler Kerne (was spezifische Verh\"altnisse von starker zu elektromagnetischer Kopplung erfordert).
    \item \textbf{Protonenstabilit\"at:} Das Proton muss innerhalb des Standardmodells stabil gegen Zerfall sein (Baryonenzahlerhaltung), sodass Wasserstoffbrennstoff \"uber kosmologische Zeitskalen verf\"ugbar bleibt ($\tau_p \gg t_\star$).
    \item \textbf{Chemische Komplexit\"at:} Die M\"oglichkeit stabiler Elektronenorbitale zur Bildung atomarer und molekularer Netzwerke.
\end{enumerate}

\subsection{Die Diproton-Katastrophe}

Detaillierte Analysen zeigen extreme Parametersensitivit\"at. Wie von Adams (2019) demonstriert, unterliegt die starke Kopplungskonstante ($\alpha_s$) einem hochsensitiven Fine-Tuning \citep{Adams2019}:

\begin{itemize}
    \item W\"are die starke Kraft um \textbf{$\sim$10\% reduziert}, w\"urde der Deuteriumkern ungebunden, was den prim\"aren nukleosynthetischen Pfad f\"ur Sternfusion vollst\"andig kappen w\"urde.
    \item W\"are die starke Kraft um \textbf{$\sim$10\% erh\"oht}, w\"urden Diprotonen (Helium-2) gebunden und hochstabil. Sterne w\"urden die tr\"age schwache Wechselwirkung vollst\"andig umgehen und ihren Wasserstoffbrennstoff \"uber die starke Wechselwirkung verbrennen, was Hauptreihensterne zum Explodieren br\"achte oder sie in einem Bruchteil ihrer heutigen Lebensdauer ausbrennen lie\ss{}e.
\end{itemize}

In der FST-I-Lesart ist die Vermeidung der Diproton-Katastrophe eine Einschr\"ankung, die langsame, dosierte stellare Energiefreisetzung m\"oglich macht. Sie beweist f\"ur sich genommen nicht, dass das totale integrierte Funktional $\mathcal{S}$ \"uber die Epoche $[t_{BB}, t_\Lambda]$ global maximiert wird.

\subsection{Die Triple-Alpha-Resonanz}

Die Produktion kritischer schwerer Elemente beruht vollst\"andig auf der Triple-Alpha-Reaktion \citep{Adams2019}. Die Gangbarkeit dieses Prozesses erfordert die pr\"azise Existenz eines $0^+$-Resonanzenergiezustands im Kohlenstoff-12-Kern. W\"urde $\alpha_{EM}$ so variiert, dass diese Resonanzenergie au\ss{}erhalb eines engen Fensters von $\sim 100$ keV bis $\sim 400$ keV verschoben w\"urde, w\"urde stellare Nukleosynthese den f\"ur komplexe Chemie und Staubbildung notwendigen Kohlenstoff und Sauerstoff nicht produzieren, was die Opazit\"ats- und K\"uhlmechanismen von Galaxien radikal ver\"andern w\"urde.

\subsection{Konkrete Toy-Variation: Die Elektronenmasse}

Betrachten wir eine heuristische Variation der Elektronenmasse, $m_e \to m_e(1+\epsilon)$ mit $|\epsilon| \ll 1$, w\"ahrend andere Parameter konstant gehalten werden.

\textbf{Positive Variation ($m_e \uparrow$):} Eine Erh\"ohung der Elektronenmasse verengt die Atomorbitale und ver\"andert die Opazit\"at des Sterninneren. Dies beschleunigt die Fusionsrate ($L_\star \sim \Gamma_{fusion}$), was zu drastisch k\"urzeren stellaren Lebensdauern ($t_\star$) f\"uhrt. W\"ahrend die momentane Rate $\dot{S}$ sprunghaft ansteigen k\"onnte, sinkt das integrierte Funktional $\mathcal{S}[\vec{p}]$ \"uber den kosmologischen Zeitrahmen signifikant aufgrund vorzeitigen stellaren Ausbrennens.

\textbf{Negative Variation ($m_e \downarrow$):} Eine Verringerung der Elektronenmasse vergr\"o\ss{}ert Atomradien und \emph{erh\"oht} die Thomson-Opazit\"at ($\kappa \propto m_e^{-2}$, also niedrigeres $m_e$ ergibt h\"oheres $\kappa$), was die Sternh\"ullen undurchsichtiger macht. Dies erh\"oht die Gleichgewichts-Kerntemperatur und steigert die Gamow-Tunnelrate $\Gamma_{\text{fusion}} \propto \exp(-\sqrt{E_G / k_B T_c})$, wobei die Gamow-Energie $E_G = (\pi \alpha)^2 m_r c^2$ von der reduzierten Kernmasse abh\"angt, nicht direkt von $m_e$. Die erh\"ohte Opazit\"at f\"angt Strahlung effektiver ein, reduziert die Photonen-Diffusionsrate und senkt dadurch die Nettoleuchtkraft. Die resultierenden l\"angeren Hauptreihen-Lebensdauern f\"uhren zu h\"oherer integrierter Entropieproduktion---konsistent mit dem monotonen Anstieg von $\mathcal{S}_{\text{total}}$ bei abnehmendem $m_e$ in Abschnitt~\ref{sec:stellar_model}. Das beobachtete $m_e$ wird daher nicht durch stellare Entropieproduktion allein fixiert, sondern durch mehrskalige Einschr\"ankungen aus der Atom- und Chemiephysik.

Die beobachtete Feinstrukturkonstante liegt nahe dem Optimum des integrierten Entropiefunktionals $\mathcal{S}$ bei festem $m_e/m_p$ (Abschnitt~\ref{sec:predictions}, Tabelle~\ref{tab:alpha_scan}: $\mathcal{S}(\alpha_{\text{obs}}) = 98{,}8\%$ des Maximums). F\"ur $m_e/m_p$ ist das Bild subtiler: Das Sternmodell zeigt monoton steigendes $\mathcal{S}$ bei abnehmendem $m_e$ (Abschnitt~\ref{sec:stellar_model}), sodass die beobachtete Elektronenmasse nicht durch stellare Entropie allein festgelegt wird, sondern durch mehrskalige Einschr\"ankungen aus der Atom- und Chemiephysik.

\subsection{Verbindung zum Nash-Gleichgewicht}

Falls ein solches eingeschr\"anktes Variationsextremum existiert, kann es als Nash-Gleichgewicht repr\"asentiert werden. Jeder fundamentale Parameter in $\vec{p}$ bildet auf eine ``Strategie'' im formal-spieltheoretischen Sinne ab---d.\,h.\ eine Konfigurationsvariable, deren zul\"assiger Wert durch die Stabilit\"atsbedingung eingeschr\"ankt wird.

Da die Parameter hochgekoppelt sind (z.\,B.\ erfordert stellare Fusion das Zusammenspiel der starken Kraft f\"ur die Bindung, der schwachen Kraft f\"ur die Proton-zu-Neutron-Umwandlung und des Elektromagnetismus f\"ur Strahlung), w\"are ein echtes Extremum von $\mathcal{S}$ ein kooperativer Zustand. Der nicht-triviale Gehalt dieser Nash-Charakterisierung liegt nicht in der Payoff-Definition (was zirkul\"ar w\"are; siehe den Zirkularit\"atsvorbehalt in Abschnitt~\ref{sec:mepp}), sondern in der \emph{Stabilit\"ats}-Aussage: einseitige Parametervariationen sollten irgendwann auf gekoppelte strukturelle Voraussetzungen treffen (nukleare Stabilit\"at, atomare Bindung, stellare Langlebigkeit)---eine Eigenschaft, die durch Parametersensitivit\"atsstudien getestet werden muss \citep{Adams2019}.

Dieser Rahmen reformuliert statisches, anthropisches Fine-Tuning als testbare Frage nach dynamischer, kooperativer Stabilit\"at.

\begin{assumption}[Potentialspiel-Struktur]\label{ass:potential-game}
Das thermodynamische Spiel, definiert durch die Sektor-Auszahlungen $u_i(p) = \dot{S}_i(p)$, besitzt eine \emph{Potentialfunktion} $S(p)$ mit
\[
  u_i(p) = \frac{\partial S}{\partial p_i}(p) \quad \forall\, i.
\]
Unter dieser Annahme reduziert sich die Best-Response-Dynamik auf Gradientenaufstieg in~$S$, der Hessian $H_{ij} = \partial_{p_i}\partial_{p_j} S$ dient als Jacobi-Matrix der Dynamik, und Resolvent-Stabilit\"at ist wohldefiniert. Ohne diese Annahme ist $H$ kein physikalisch sinnvoller Stabilit\"atsoperator.
\end{assumption}

\begin{remark}[Physikalische Motivation f\"ur die Potentialstruktur]
\label{rem:potential-motivation}
Nahe dem thermodynamischen Gleichgewicht garantieren die Onsager-Reziprozit\"atsbeziehungen, dass die Matrix der Transportkoeffizienten symmetrisch ist, $L_{ij} = L_{ji}$, was impliziert, dass die Entropieproduktion $\dot{S} = \sum_{ij} L_{ij} X_i X_j$ aus einem quadratischen Potential in den thermodynamischen Kr\"aften $X_i$ stammt \citep{Prigogine1967}. Die Potentialspiel-Annahme erweitert diese Struktur auf das fern vom Gleichgewicht liegende Regime, in dem Onsager-Reziprozit\"at im Allgemeinen nicht mehr gilt und die Existenz einer globalen Potentialfunktion nicht garantiert ist \citep{GrinsteinLinsker2007}. Diese Erweiterung ist die zentrale ungetestete strukturelle Annahme des Rahmens. Empirische Plausibilit\"at ergibt sich aus der Beobachtung, dass in bekannten physikalischen Systemen mit hoher Symmetrie (z.\,B.\ der Eichstruktur des Standardmodells) die Kopplungsmatrix zwischen Sektoren N\"aherungssymmetrien aufweist, die eine Potentialzerlegung plausibel machen---eine rigorose Herleitung f\"ur den vollen Standardmodell-Parameterraum bleibt jedoch offen.
\end{remark}

\begin{remark}[Fine-Tuning als Resolventenstabilit\"at]
\label{rem:fine-tuning-resolvent}
Die Resolventenperspektive \citep{GeigerRH2026b} liefert ein strukturelles Vokabular f\"ur Fine-Tuning, das weder teleologisch noch anthropisch ist. In f\"uhrender Ordnung k\"onnen kleine Parametervariationen nahezu neutral sein---viele nahe Konfigurationen erzeugen vergleichbare Entropie. Im konjekturalen FST-I-Mechanismus w\"urden diese Variationen in zweiter Ordnung durch Resonanzkopplungen zwischen Parametersektoren unterscheidbar (analog zur spektralen Aufspaltung zweiter Ordnung in der Analyse der Weil-quadratischen Form, wo eine Stabilit\"atsl\"ucke durch Kopplung zwischen komplement\"aren Unterr\"aumen entsteht). Fine-Tuning w\"are dann eine m\"ogliche Signatur von Resolventenstabilit\"at: Die beobachteten Parameter w\"urden einen eingeschr\"ankten Fixpunkt besetzen, an dem Kreuzkopplungen zweiter Ordnung zwischen entarteten Moden gleichzeitig stabilisiert sind. Dies bleibt Hypothese, bis projizierter Gradient, Tangential-Hessian und KKT-Bedingungen gepr\"uft sind.
\end{remark}

\subsection{G\"ultigkeitsbereich und Einschr\"ankungen}

Wir beanspruchen keine vollst\"andige Herleitung des Standardmodells aus ersten Prinzipien oder eine A-priori-Berechnung des exakten Werts von $\alpha_{EM}$. Vielmehr testen wir, ob die beobachteten Parameterwerte mit einer eingeschr\"ankten Hochentropie-Stabilit\"atsregion kompatibel sein k\"onnen.

\textbf{Kritische Einschr\"ankung:} W\"ahrend Abschnitt~\ref{sec:stellar_model} eine erste semi-analytische Berechnung liefert, die zeigt, dass $\mathcal{S}_{\text{total}}(\alpha_{\text{obs}})$ 98,8\% des uneingeschr\"ankten stellaren Maximums erreicht, adressiert diese Berechnung nur einen einzelnen Entropiekanal (Hauptreihen-Sternstrahlung) und variiert nur zwei Parameter ($\alpha$ und $m_e/m_p$). Ein vollst\"andiger Test der FST-I-Hypothese erfordert die Evaluierung von $\mathcal{S}[\vec{p}] = \int_{t_{BB}}^{t_\Lambda} \dot{S}(\vec{p}, t)\,dt$ \"uber alle 19 Standardmodell-Parameter und alle Entropieproduktionskan\"ale (stellar, gravitationell, chemisch, nuklear). Eine solche Berechnung w\"urde umfassen: (i) Sternentwicklungscodes mit nicht-standardm\"a\ss{}igen nuklearen Parametern, (ii) Galaxienbildungssimulationen und (iii) Integration \"uber die gesamte kosmologische Geschichte.

Das semi-analytische Ergebnis ist ermutigend---das beobachtete $\alpha$ ist nahezu optimal f\"ur stellare Entropie---stellt aber noch keine vollst\"andige Demonstration dar. Insbesondere zeigt die monotone $m_e/m_p$-Abh\"angigkeit, dass mehrskalige Einschr\"ankungen (Atomstabilit\"at, Chemie-Gangbarkeit) ber\"ucksichtigt werden m\"ussen. Die Erweiterung dieser Berechnung ist die prim\"are Aufgabe f\"ur zuk\"unftige Arbeit.

FST-I ist ein funktionaler Rahmen. Er diktiert nicht die fundamentale mikroskopische Topologie irgendeiner tieferen Theorie. Stattdessen schl\"agt er einen Randbedingungs-Test vor: Effektive Niederenergie-Parameter, die aus einer tieferen Theorie hervorgehen, sollten mit eingeschr\"ankter Entropieproduktions-Stabilit\"at \"uber kosmologische Zeitskalen kompatibel sein.

% ===================================================================
\section{Symmetriebrechung als Phasen\"ubergang im Spiel}
\label{sec:symmetry}
% ===================================================================

Abschnitt~\ref{sec:quantum} formulierte ausgew\"ahlte Quantenstrukturen in der Sprache gemischter Strategien um, und Abschnitt~\ref{sec:finetuning} definierte ein Kandidatenfunktional f\"ur die Standardmodell-Parameter. Dieser Abschnitt behandelt den dynamischen Nexus zwischen beiden: den Ursprung von Masse und Struktur.

\subsection{Symmetrie als entartetes Gleichgewicht}

Im spieltheoretischen Rahmen entspricht eine symmetrische Phase einem entarteten Nash-Gleichgewicht: mehrere Strategien ergeben identische thermodynamische Payoffs. Die ungebrochene Symmetrie repr\"asentiert ein flaches Payoff-Plateau im Konfigurationsraum. Da keine einzelne Strategie thermodynamisch bevorzugt ist, ist das System marginal stabil.

W\"ahrend der fr\"uhesten Epochen des Universums, vor der elektroschwachen Epoche, \"ubertraf die Umgebungsw\"armeenergiedichte bei weitem die Ruhemasseenergien aller wechselwirkenden Teilchen \citep{HiggsVEV}. In diesem Hochtemperaturregime wird ein masseloses, hochsymmetrisches Plasma hier als Kandidat einer Hochdurchsatz-Konfiguration f\"ur schnelle Energiedissipation gelesen, da es im thermischen Standardbild die verf\"ugbaren Freiheitsgrade und Kollisionsraten maximiert \citep{HiggsVEV2024}.

\subsection{Elektroschwache Symmetriebrechung als Instabilit\"at}

Wenn sich die \"au\ss{}eren kosmologischen Randbedingungen \"andern---durch die Expansion und Abk\"uhlung des Universums---wird das zuvor symmetrische Gleichgewicht instabil. Oberhalb der kritischen Temperatur ($T > T_c$) liegt das Minimum des effektiven Potentials bei $\phi = 0$, entsprechend dem einzigen stabilen Gleichgewicht. Wenn die Temperatur unter $T_c$ f\"allt, geht $\phi = 0$ von einem stabilen Minimum zu einem lokalen Maximum (instabiles Gleichgewicht) \"uber, und neue Minima erscheinen bei $\phi \neq 0$.

Das selbstwechselwirkende Skalarfeldpotential, das diesen \"Ubergang bestimmt, ist:\footnote{Wir folgen der Peskin-\&-Schroeder-Konvention \citep{HiggsVEV}, bei der $V(\phi) = -\mu^2 |\phi|^2 + \lambda |\phi|^4$ mit $\mu^2 > 0$ f\"ur die gebrochene Phase gilt. In der hier verwendeten reellen Skalar-Kurzform sind die Faktoren $\frac{1}{2}$ und $\frac{1}{4}$ konventionelle Normierungen.}
\begin{equation}
V(\phi) = -\frac{1}{2}\mu^2 \phi^2 + \frac{1}{4}\lambda \phi^4
\end{equation}

Wenn das Universum abk\"uhlt, entwickelt sich der Massenparameter $\mu^2$, durchl\"auft Null und wird positiv (in der Konvention, bei der das negative Vorzeichen explizit ist). Dies provoziert spontane Symmetriebrechung \citep{HiggsVEV2024}. Im Kontext des thermodynamischen Spiels wird der \"Ubergang des Vakuums in eine gebrochene Symmetriephase als Stabilit\"ats\"ubergang gelesen, der mit Entropieproduktions-Constraints kompatibel ist, nicht als aus MEPP allein abgeleitete Konsequenz.

\subsection{Das Higgs-VEV als Attraktor des Spiels}

Der Higgs-Vakuumerwartungswert (VEV) wird hier als Kandidat f\"ur einen stabilen Fixpunkt des Level-1-Spiels interpretiert. Die physikalische Aussage bleibt die etablierte: Ein von Null verschiedener Erwartungswert, $\langle\phi\rangle \neq 0$, minimiert das Higgs-Potential in der gebrochenen Phase. Die zus\"atzliche FST-I-Behauptung ist nur, dass dieser Fixpunkt m\"oglicherweise als Teil eines eingeschr\"ankten thermodynamischen Stabilit\"atsmusters lesbar ist.

Der exakte numerische Wert des Higgs-VEV, ungef\"ahr 246 GeV, dient als fundamentaler Massengenerator f\"ur die W- und Z-Bosonen sowie die Fermionen \"uber Yukawa-Kopplungen \citep{HiggsVEV}. Dieser pr\"azise Skalierungsparameter trennt das masselose Photon von den massiven Mediatoren der schwachen Kraft, wodurch effektiv die Disparit\"at in den Kraftst\"arken erzeugt wird, die langfristige strukturelle Stabilit\"at der Kernmaterie erm\"oglicht.

In der FST-I-Lesart ist Massenerzeugung thermodynamisch relevant, weil sie Entropieproduktion h\"oherer Ordnung erm\"oglicht. Ein ausschlie\ss{}lich von masselosen Teilchen bev\"olkertes Universum w\"urde auf unbestimmte Zeit strahlungsdominiert bleiben und den gravitativen Kollaps zu gebundenen Strukturen (Sternen, Galaxien) verhindern, wodurch nachhaltige lokalisierte Dissipation nicht verf\"ugbar w\"are.

\subsection{Verbindung zur kosmologischen S\"attigung}

Diese Dynamik hat eine formale Analogie zur Dynamik auf kosmologischer Skala. J\"ungste kosmologische Rekonstruktionen, die den kinematischen Jerk-Parameter ($j = 1$ f\"ur $\Lambda$CDM) mit Skalarfeld-Dynamik verkn\"upfen, diskutieren eine theoretische kosmische Variation im Proton-zu-Elektron-Massenverh\"altnis ($\mu = m_p/m_e$), die durch die Evolution des Higgs-VEV \"uber die Zeit ($v(z)$) angetrieben wird \citep{HiggsVEV2024}. Theoretische Berechnungen dieser Variation passen innerhalb der Beobachtungseinschr\"ankungen aus Quasar-Absorptionsspektren und liefern einen Hinweis, dass der mikroskopische Symmetriebrechungs-Ordnungsparameter ($v$) dynamisch an die makroskopische Expansionsrate gekoppelt bleiben k\"onnte.

Im Kr\"ummungs-R\"uckgabemodell (CRM \cite{GeigerCRM2026}, siehe auch \cite{FST-Unified2026}) ergibt die gr\"obste Skala (``Level~0'') der Raum-Zeit-Ausrichtung eine $\tanh$-S\"attigungskurve als ihre Mean-Field-L\"osung. Auf der n\"achstfeineren Skala (``Level~1'') dient das quartische Higgs-Potential als effektive Beschreibung. FST-I behandelt beide Ph\"anomene als m\"oglicherweise verwandte Stabilit\"atsmuster auf verschiedenen Skalen, nicht als bewiesene Ableitungen aus derselben Optimierungslogik.

\subsection[Kosmologische Phasen\"uberg\"ange als sukzessive Nash-Gleichgewichte]{Kosmologische Phasen\"uberg\"ange als sukzessive\\ Nash-Gleichgewichte}

Jeder kosmologische Phasen\"ubergang wird in diesem Vokabular als m\"oglicher \"Ubergang zwischen Nash-Gleichgewichten gelesen:

\begin{center}
\begin{tabularx}{\textwidth}{@{}>{\raggedright\arraybackslash}p{3.4cm}YY@{}}
\toprule
\textbf{\"Ubergang} & \textbf{Strategischer Wechsel} & \textbf{Konsequenz} \\
\midrule
Elektroschwach & Relaxiert in VEV-Konfiguration & Erm\"oglicht gravitative Klumpenbildung \\
QCD-Confinement & Relaxiert in Bindungszustandskonfiguration & Erzeugt stabile Protonen-Reservoirs \\
Rekombination & Relaxiert in neutrale-Atom-Konfiguration & Erzeugt CMB-W\"armesenke \\
\bottomrule
\end{tabularx}
\end{center}

Jeder Phasen\"ubergang repr\"asentiert eine Reduktion reiner strategischer Freiheit (Symmetrieerniedrigung) im Austausch gegen eine Erh\"ohung irreversibler, robuster Entropieproduktion.

% ===================================================================
\section{Die Rolle von Information und Dekoh\"arenz}
\label{sec:decoherence}
% ===================================================================

W\"ahrend Abschnitt~\ref{sec:quantum} die kommutative Gemischte-Strategie-Analogie f\"ur Quantenzust\"ande eingef\"uhrt hat, ist es notwendig, die pr\"azise informationale Mechanik zu untersuchen, die die Strategieselektion antreibt. In der Standard-Quantentheorie wirkt die Umgebung als Reservoir und Monitoring-Kanal; in FST-I wird dies analog als kompetitive Informationsgrenze in einem Nullsummenspiel modelliert \citep{ValentiniWestman2005}.

\subsection{Dekoh\"arenz als strategische Durchsetzung}

Die Dekoh\"arenztheorie liefert den physikalischen Mechanismus, der die spieltheoretische Analogie motiviert. W\"ahrend ein Quantensystem evolviert, interagiert es kontinuierlich mit dem umgebenden Vakuum und den Thermalb\"adern. Diese Interaktion verschr\"ankt den Zustand des Systems mit den unbeobachtbaren Freiheitsgraden der Umgebung \citep{MeasurementClassical}.

Da dem lokalen Beobachter der Zugang zur gesamten universellen Wellenfunktion verwehrt ist, verliert die reduzierte Dichtematrix des Systems rasch ihre Nichtdiagonal-Interferenzterme durch Ausspuren \"uber Umgebungszust\"ande \citep{Lindblad1976}. F\"ur mesoskopische Systeme reichen Dekoh\"arenzzeitskalen von ${\sim}10^{-23}\,$s (stark wechselwirkende Hadronen in einem thermischen Bad) bis ${\sim}10^{-17}\,$s (molekulare Superpositionen bei Raumtemperatur), was sicherstellt, dass die Strategieselektion auf Zeitskalen stattfindet, die weit k\"urzer sind als jede makroskopische Beobachtung.

\subsection{Einselection als Nash-Durchsetzung}

Dieser Prozess, genannt Einselection (umgebungsinduzierte Superselektion), ist nicht nur passiver Zerfall; im hier verwendeten Vokabular wird er als Durchsetzung stabiler Zeigerbasis-Struktur modelliert \citep{Zurek2003}. Die Umgebung \"uberwacht selektiv bestimmte Observablen (typischerweise solche, die mit dem Wechselwirkungs-Hamiltonian kommutieren, wie Position), und unterdr\"uckt Superpositionen, die nicht mit stabilen ``Zeigerzust\"anden'' \"ubereinstimmen.

Die Realisierung eines spezifischen Messergebnisses ist daher die thermodynamische Kulmination des Spiels, wobei das informationale Potential einer ausgedehnten, nicht-kollabierten gemischten Strategie irreversibel in einen singul\"aren physikalischen Zustand umgewandelt wird, was eine lokalisierte Spitze in der Entropie erzeugt und den thermodynamischen Payoff sichert \citep{Lindblad1976}.

% ===================================================================
\section{Verbindungen zu bestehenden Ans\"atzen}
\label{sec:connections}
% ===================================================================

\subsection{Entropische Gravitation: Verlinde und Padmanabhan}

Verlindes Vorschlag der entropischen Gravitation \citep{Verlinde2011} legt nahe, dass Gravitationskr\"afte aus entropischen Gradienten entstehen. Padmanabhans Programm \citep{Padmanabhan2010} leitet gravitationelle Feldgleichungen aus thermodynamischen Identit\"aten ab. Unser Rahmen nutzt diese Logik als Analogie f\"ur den Standardmodell-Parameterraum, nicht als etablierte Ableitung.

\subsection{Umgebungsinduzierte Superselektion: Zurek}

Zureks Dekoh\"arenzprogramm \citep{Zurek2003} erkl\"art klassische Emergenz durch Umgebungs\"uberwachung. Unser Rahmen liest Zeigerzust\"ande analog als Nash-stabile Strategien, wobei Einselection als strukturelles Gegenst\"uck zur Unterdr\"uckung instabiler Strategien dient.

\subsection{Free Energy Principle: Friston}

Fristons Free Energy Principle \citep{Friston2010} schl\"agt vor, dass biologische Systeme ihre variationelle freie Energie minimieren. Wir behandeln dies als biologisches Vergleichsmodell mit verwandter Variationsstruktur, nicht als direkte Grobk\"ornungsfolge der Teilchenebene.

\paragraph{Active Inference als stochastisches optimales Kontrollproblem.}
Active Inference l\"asst sich als stochastisches optimales Kontrollproblem in den MFG-Rahmen einbetten \citep{LasryLions2006,LasryLions2007}. Jeder Agent h\"alt interne \"Uberzeugungen $\mu_t \in \mathbb{R}^n$ (hinreichende Statistiken von $q(s|\mu)$) und w\"ahlt Aktionen $a_t \in \mathcal{U}$, die Blanket-Zust\"ande modifizieren. Die Dynamik folgt
\[
  d\mu_t = f(\mu_t, a_t, m_t)\,dt + \sigma\,dW_t,
\]
wobei $m_t$ die Mean-Field-Verteilung \"uber die Zust\"ande der anderen Agenten ist. Das Kostenfunktional ist die variationelle freie Energie:
\[
  L(\mu, a, m) = \mathrm{KL}\bigl(q(s|\mu) \,\|\, p(s|o,a,m)\bigr) - \ln p(o|a,m).
\]
Die Hamilton--Jacobi--Bellman-Gleichung f\"ur die Wertfunktion $V(\mu,t) := \inf_{a} \mathbb{E}\bigl[\int_0^T L\,dt + \Phi(\mu_T)\bigr]$ lautet
\[
  -\partial_t V = \inf_a\bigl[L(\mu,a,m) + \nabla_\mu V \cdot f(\mu,a,m)\bigr] + \tfrac{\sigma^2}{2}\Delta_\mu V.
\]
Gekoppelt mit der Fokker--Planck-Gleichung f\"ur $m$ ergibt sich ein MFG-System, dessen Fixpunkt $(V^*, m^*)$ dem Nash-Gleichgewicht des Multi-Agenten-Systems entspricht. Wenn die freie Energie additiv zerlegbar ist, $F = \sum_i F_i$, reduziert sich das MFG auf ein Potentialspiel (Annahme~\ref{ass:potential-game}), was die formale Br\"ucke zwischen Active Inference und dem FST-Rahmen liefert.

\subsection{Fine-Tuning-Studien: Adams, Tegmark}

Adams \citep{Adams2019} und Tegmark et al.\ \citep{Tegmark2006} haben den Parameterraum m\"oglicher Universen systematisch erforscht und gangbare Regionen kartiert, in denen stabile Kerne, Atome und Sterne existieren k\"onnen. Tegmark et al.\ identifizieren 31 dimensionslose Parameter und zeigen, dass viele in engen Bereichen liegen m\"ussen, damit Strukturen entstehen. Unser Beitrag besteht darin, ein organisierendes Prinzip vorzuschlagen: Die beobachteten Parameter k\"onnten in einer eingeschr\"ankten Hochentropie-/Stabilit\"atsregion liegen; der aktuelle Alpha-Scan st\"utzt aber nur einen Pilot-Konsistenzcheck, kein globales thermodynamisches Optimum.

\subsection{Alternative Rahmungen und ihre Beziehung zu FST-I}

Mehrere alternative Erkl\"arungen f\"ur Fine-Tuning m\"ussen ber\"ucksichtigt und mit FST-I verglichen werden:

\begin{itemize}
    \item \textbf{Minimale Entropieproduktion (MinEP):} Prigogines Theorem der minimalen Entropieproduktion \citep{Prigogine1967} gilt f\"ur Systeme nahe dem thermodynamischen Gleichgewicht. Es sagt das Gegenteil von MEPP f\"ur gleichgewichtsnahe Zust\"ande vorher. FST-I behauptet, dass die kosmologische Epoche der Strukturbildung weit vom Gleichgewicht entfernt ist und sich daher im MEPP-Regime befindet; diese Behauptung erfordert Verifikation.
    \item \textbf{Anthropisches Prinzip:} Weinberg \citep{Weinberg1987} und andere argumentieren, dass beobachtete Parameter durch Beobachterexistenz selektiert werden, nicht durch Optimierung. FST-I testet die schärfere Hypothese, dass beobachtete Parameter in eingeschränkten Hochentropie-/Stabilitätsregionen liegen und nicht bloß irgendwo im anthropischen Fenster. Der entscheidende empirische Unterschied besteht darin, ob die beobachteten Parameter an der \emph{Grenze} des gangbaren Fensters liegen (konsistent mit anthropischer Selektion) oder an einem unabhängig eingeschränkten inneren Stabilitätspunkt von $\mathcal{S}[\vec{p}]$ (konsistent mit FST-I). Dies zu testen erfordert die numerischen Berechnungen, die in Abschnitt~\ref{sec:predictions} beschrieben werden.
    \item \textbf{String-Landscape:} Susskind \citep{Susskind2005} schlägt vor, dass die Landschaft möglicher String-Vakua, kombiniert mit ewiger Inflation, eine statistische Erklärung für Fine-Tuning liefert. FST-I erfordert keine Stringtheorie und formuliert ein eigenständiges vorläufiges Prüfkriterium: Nachdem unabhängige Gangbarkeits-Constraints fixiert sind, sollten Entropieproduktion und Stabilität in der Nähe der beobachteten Werte besser ausfallen als an passenden Kontrollpunkten.
    \item \textbf{Kosmologische natürliche Selektion:} Smolin \citep{Smolin1992} schlägt vor, dass sich Universen über Schwarze Löcher reproduzieren, wobei Parameter selektiert werden, die die Schwarzloch-Produktion maximieren. Dies liefert zwar einen evolutionären Selektionsmechanismus, leitet jedoch keine dynamischen Gleichungen ab und sagt keine spezifischen Parameterwerte vorher. FST-I ersetzt reproduktive Fitness durch thermodynamische Stabilität als vorgeschlagenes Selektionskriterium und liefert rechnerische Diagnostiken statt etablierter Parametervorhersagen.
    \item \textbf{Zellulärer-Automat-Interpretation:} 't~Hooft \citep{tHooft2016} schlägt vor, dass die Quantenmechanik aus einer zugrunde liegenden deterministischen Zellulären-Automat-Dynamik hervorgeht. Während FST-I keine versteckten Variablen einführt, suchen beide Programme ein tieferes Organisationsprinzip hinter der axiomatischen Struktur der Quantenmechanik. Der wesentliche Unterschied besteht darin, dass FST-I die Quantenstruktur durch Optimierungs- und Stabilitätsbedingungen liest, ohne mikroskopischen Determinismus zu erfordern.
\end{itemize}

Tabelle~\ref{tab:framework-comparison} bietet einen kompakten Vergleich von FST-I mit den wichtigsten alternativen Rahmenwerken, die Fine-Tuning oder den Ursprung physikalischer Gesetze adressieren.

\begin{table}[ht]
\centering
\small
\caption{Vergleich von Rahmenwerken zur Erkl\"arung von Fine-Tuning und dem Ursprung physikalischer Parameter.}
\label{tab:framework-comparison}
\begin{tabularx}{\textwidth}{@{}>{\raggedright\arraybackslash}p{3.0cm}YYY@{}}
\toprule
\textbf{Rahmenwerk} & \textbf{Mechanismus} & \textbf{Befund / Teststatus} & \textbf{Status} \\
\midrule
Anthropisches Prinzip \citep{Weinberg1987} & Beobachterselektion auf $\Lambda$ & $\Lambda$-Schranke; keine Dynamik & Kein falsifizierbarer Mechanismus \\[4pt]
Kosmologische nat\"urliche Selektion \citep{Smolin1992} & Schwarzloch-Reproduktion selektiert Parameter & Maximiert Schwarzloch-Anzahl & Nicht falsifiziert; keine Dynamik abgeleitet \\[4pt]
Free Energy Principle \citep{Friston2010} & Variationelle Inferenz minimiert \"Uberraschung & Verhalten auf Agentenebene & Nicht auf fundamentale Physik angewandt \\[4pt]
Zellul\"arer Automat QM \citep{tHooft2016} & Deterministische Substrate f\"ur QM & Diskretheitssignaturen & Keine Parameterselektion \\[4pt]
\textbf{FST-I (diese Arbeit)} & Stabilit\"atsbeschr\"ankte Entropiemaximierung & $\alpha$ im breiten Entropieplateau (98,8\,\%); $m_e/m_p$ erfordert mehrskalige Constraints (47\,\% in 2D) & Falsifizierbar \"uber Mehrparameter-St\"orung; nur Pilotstudie \\
\bottomrule
\end{tabularx}
\end{table}

FST-I kombiniert (i)~einen konkreten Stabilit\"atsmechanismus (Nash-Gleichgewicht eines thermodynamischen Spiels) mit (ii)~einem ersten Pilotstudien-Konsistenzcheck (das $\alpha$-Scan-Plateau bei 98,8\,\% Entropieoptimalit\"at f\"ur einen Einzelparameter in einem vereinfachten Sternmodell) und (iii)~einem expliziten Falsifikationskriterium (Abschnitt~\ref{sec:predictions}). Wir betonen, dass (ii) keine Vorhersage ist, sondern ein post-hoc-Konsistenzcheck mit bekannter Physik; der Gamow-Tunnelfaktor, der das 98,8\,\%-Resultat tr\"agt, ist selbst Standard-Kernphysik. Das anthropische Prinzip und die kosmologische nat\"urliche Selektion liefern Selektionsdruck, aber keine dynamischen Gleichungen; das Free Energy Principle operiert auf der biologischen Skala und wurde nicht auf fundamentale Konstanten erweitert; 't~Hoofts Zellul\"arer-Automat-Programm adressiert den Ursprung der Quantenmechanik, aber nicht die Werte der Standardmodell-Parameter.

Der FST-I-Rahmen kann prinzipiell von diesen Alternativen durch die geplanten numerischen Tests in Abschnitt~\ref{sec:predictions} unterschieden werden. Allerdings wurde diese Unterscheidung noch nicht rechnerisch demonstriert.

\begin{remark}[Nash-Gleichgewichte als stabile Fixpunkte zweiter Ordnung]
\label{rem:degenerate-perturbation}
Die Entropieproduktionslandschaft ist in f\"uhrender Ordnung flach (Tabelle~\ref{tab:alpha_scan}: $\mathcal{S}/\mathcal{S}_{\max} > 0{,}92$ \"uber den gangbaren Bereich). Die beobachtete Konfiguration k\"onnte durch Kreuzkopplungen zweiter Ordnung zwischen Parametersektoren selektiert werden, analog zur entarteten St\"orungstheorie. Diese Konjektur erfordert die in der Schlussfolgerung genannte Mehrparameter-Hesse-Berechnung; sie wurde bislang nicht verifiziert.
\end{remark}

% ===================================================================
\section{Testbare Vorhersagen und Falsifizierbarkeit}
\label{sec:predictions}
% ===================================================================

\subsection{Hauptvorhersage: Entropiedominanz des Standardmodells}

\begin{conjecture}[Entropiedominanz]
\label{conj:entropy-dominance}
Der beobachtete Standardmodell-Parametersatz $\vec{p}_{\mathrm{SM}}$ entspricht einem lokalen Maximum des integrierten Entropieproduktions-Funktionals $\mathcal{S}$ innerhalb der strukturell stabilen Region des Parameterraums:
\begin{equation}
\mathcal{S}(\mathrm{SM}) > \mathcal{S}(\mathrm{SM}')
\end{equation}
f\"ur jeden variierten Parametersatz $\mathrm{SM}'$, wobei $|\Delta p_i| \geq 10\%$ f\"ur mindestens einen Parameter $p_i$, unter strukturellen Stabilit\"atseinschr\"ankungen (nukleare Stabilit\"at, atomare Bindung, chemische Komplexit\"at).
\end{conjecture}

\subsection{Konkrete numerische Tests}

\textbf{P1a: Elektron-zu-Proton-Massenverh\"altnis ($m_e/m_p$):} Das Durchlaufen von Sternentwicklungscodes \"uber ein Gitter von Werten sollte zeigen, ob integriertes $\Delta S$ nahe dem beobachteten Wert durch mehrskalige Kompatibilit\"at eingeschr\"ankt wird: stellare Entropieproduktion ist in $m_e/m_p$ monoton (Abschnitt~\ref{sec:stellar_model}), daher muss das beobachtete Verh\"altnis, wenn \"uberhaupt, durch atomare und chemische Stabilit\"atsgrenzen statt durch einen einzelskaligen Entropiepeak erkl\"art werden. Signifikante Abweichungen werden voraussichtlich die stellare Energieerzeugung oder die Chemie in ineffiziente beziehungsweise instabile Regime zwingen \citep{Adams2019}.

\textbf{P1b: Feinstrukturkonstante ($\alpha_{EM}$):} Wir erwarten ein breites Optimierungsplateau nahe dem beobachteten Wert. Tabelle~\ref{tab:alpha_scan} ist im Pilotmodell konsistent mit $\mathcal{S}/\mathcal{S}_{\max} > 0{,}92$ f\"ur $0{,}5 \leq \alpha/\alpha_{\text{obs}} \leq 2{,}0$. Ein signifikanter Abfall ist erst jenseits von $\alpha/\alpha_{\text{obs}} \lesssim 0{,}5$ oder $\gtrsim 3{,}0$ zu erwarten, und zwar aufgrund des Verlusts stabiler chemischer Bindung, atomarer Stabilit\"at oder verk\"urzter stellarer Lebensdauern.

\subsection{Semi-analytisches Sternmodell: Quantitative Ergebnisse}
\label{sec:stellar_model}

Als \emph{Pilotstudie} f\"ur einen ersten quantitativen Test der Entropiedominanz-Vorhersage konstruieren wir ein semi-analytisches Sternmodell, das die integrierte Entropieproduktion $\mathcal{S}_{\text{total}} = \int_0^{t_{\text{MS}}} \dot{S}(t)\,dt \approx E_{\text{nuc}} / T_{\text{eff}}$ als Funktion der Feinstrukturkonstante $\alpha$ und des Elektron-zu-Proton-Massenverh\"altnisses $m_e/m_p$ berechnet. Die N\"aherung $\mathcal{S}_{\text{total}} \approx E_{\text{nuc}} / T_{\text{eff}}$ ist g\"ultig, da die effektive Temperatur eines Hauptreihensterns sich \"uber seine Wasserstoffbrennzeit um weniger als einen Faktor zwei \"andert, sodass $T_{\text{eff}}$ f\"ur diese Absch\"atzung als n\"aherungsweise konstant behandelt werden kann.

\textbf{Modellannahmen.} Wir betrachten einen Stern mit Sonnenmasse auf der Hauptreihe, modelliert als Polytrop mit Index $n = 3$. Die Kerntemperatur wird selbstkonsistent aus dem Gleichgewicht zwischen Kernenergieerzeugung (dominiert von der pp-Kette mit Gamow-Tunnelfaktor) und radiativem Energietransport (Thomson-Opazit\"at $\kappa \propto m_e^{-2}$) bestimmt. Die Gamow-Energie f\"ur Proton-Proton-Fusion ist:
\begin{equation}
E_G = 2 m_r c^2 (\pi \alpha)^2 = (\pi \alpha)^2 m_p c^2 \approx 493\,\text{keV}
\end{equation}
wobei $m_r = m_p/2$ die reduzierte Masse des Proton-Proton-Systems ist, und der Tunnelparameter bei solarer Kerntemperatur $\tau_\odot = 3(E_G / 4 k_B T_c)^{1/3} \approx 13{,}5$ betr\"agt. Die Leuchtkraft skaliert als $L \propto T_c^{\nu}$ wobei $\nu = \tau/3 - 2/3 \approx 3{,}8$ f\"ur pp-Ketten-Bedingungen. Die Hauptreihen-Lebensdauer folgt aus $t_{\text{MS}} = E_{\text{nuc}} / L$.

\textbf{Hinweis zu absoluten Zeitskalen.} Das Modell berechnet $t_{\text{MS}}$ aus dem gesamten nuklearen Energiereservoir geteilt durch die Leuchtkraft, was Werte ergibt, die von detaillierten Sternentwicklungscodes (die ${\sim}10\,$Gyr f\"ur einen Stern mit Sonnenmasse ergeben) um einen Faktor ${\sim}7$ abweichen. Dieser systematische Offset resultiert aus der vereinfachten polytropischen Behandlung (keine Opazit\"atstabellen, keine konvektive H\"ulle, keine Metallizit\"atseffekte). Da alle Eintr\"age in Tabelle~\ref{tab:alpha_scan} mit demselben Modell berechnet werden, sind die \emph{relativen} Trends---insbesondere das Verh\"altnis $\mathcal{S}/\mathcal{S}_{\text{max}}$---robust, w\"ahrend die absoluten Zeitskalen nicht direkt mit beobachteten stellaren Lebensdauern verglichen werden sollten.

\textbf{Alpha-Variation.} Tabelle~\ref{tab:alpha_scan} fasst die Ergebnisse der Variation von $\alpha / \alpha_{\text{obs}}$ von 0,25 bis 5,0 zusammen, w\"ahrend $m_e/m_p$ und $\alpha_s$ konstant gehalten werden. Die starke Kopplungskonstante wird konstant gehalten, da ihre Variation die Kernphysik (Bindungsenergien, Fusionspfade) fundamental ver\"andern w\"urde, was vom polytropischen Sternmodell nicht erfasst werden kann. Das Schl\"usselergebnis ist, dass die integrierte Entropieproduktion $\mathcal{S}_{\text{total}}$ ein breites Maximum nahe $\alpha / \alpha_{\text{obs}} \approx 1{,}6$ aufweist, wobei der beobachtete Wert $\mathcal{S}(\alpha_{\text{obs}}) / \mathcal{S}_{\text{max}} = 0{,}988$ erreicht.

\begin{table}[ht]
\centering
\caption{Integrierte stellare Entropieproduktion vs.\ Feinstrukturkonstante. Der beobachtete Wert $\alpha_{\text{obs}} = 1/137$ erreicht 98,8\% des uneingeschr\"ankten Maximums.}
\label{tab:alpha_scan}
\begin{tabular}{cccccc}
\toprule
$\alpha / \alpha_{\text{obs}}$ & $L / L_\odot$ & $T_{\text{eff}}$ [K] & $t_{\text{MS}}$ [Gyr] & $\mathcal{S}_{\text{total}}$ [J/K] & $\mathcal{S} / \mathcal{S}_{\text{max}}$ \\
\midrule
0{,}25 & 15{,}52 & 7026 & 4{,}67 & $1{,}247 \times 10^{41}$ & 0{,}811 \\
0{,}50 & 4{,}00 & 6171 & 18{,}1 & $1{,}420 \times 10^{41}$ & 0{,}924 \\
0{,}75 & 1{,}77 & 5889 & 41{,}0 & $1{,}487 \times 10^{41}$ & 0{,}968 \\
\textbf{1{,}00} & \textbf{1{,}00} & \textbf{5772} & \textbf{72{,}5} & $\mathbf{1{,}518 \times 10^{41}}$ & \textbf{0{,}988} \\
1{,}25 & 0{,}63 & 5720 & 114 & $1{,}532 \times 10^{41}$ & 0{,}997 \\
1{,}50 & 0{,}44 & 5701 & 165 & $1{,}536 \times 10^{41}$ & 1{,}000 \\
2{,}00 & 0{,}25 & 5712 & 290 & $1{,}534 \times 10^{41}$ & 0{,}998 \\
3{,}00 & 0{,}11 & 5804 & 640 & $1{,}508 \times 10^{41}$ & 0{,}981 \\
5{,}00 & 0{,}04 & 6043 & 1813 & $1{,}450 \times 10^{41}$ & 0{,}943 \\
\bottomrule
\end{tabular}
\end{table}

\textbf{Strukturelle Stabilit\"atseinschr\"ankungen.} Das uneingeschr\"ankte Maximum bei $\alpha / \alpha_{\text{obs}} \approx 1{,}6$ muss gegen physikalische Gangbarkeit bewertet werden: f\"ur $\alpha \gtrsim 1/85$ h\"oren stabile Atome schwerer als Eisen auf zu existieren (die Feinstrukturentwicklung bricht zusammen), und f\"ur $\alpha \gtrsim 1/60$ verschiebt sich die Triple-Alpha-Resonanz aus ihrem gangbaren Fenster \citep{Adams2019}. Wenn diese Einschr\"ankungen angewandt werden, verengt sich das gangbare Parameterfenster auf ungef\"ahr $0{,}5 \lesssim \alpha/\alpha_{\text{obs}} \lesssim 2{,}0$, innerhalb dessen der beobachtete Wert nahe dem eingeschr\"ankten Maximum liegt.

\textbf{Elektronenmasse-Variation: mehrskalige Einschr\"ankungen.} Das Sternmodell zeigt eine monotone Zunahme von $\mathcal{S}_{\text{total}}$ mit abnehmender $m_e$, was das Fehlen einer intrinsischen unteren Grenze f\"ur $m_e$ auf der stellaren Skala widerspiegelt (Opazit\"at $\kappa \propto m_e^{-2}$ kontrolliert den Leuchtkraft-Lebensdauer-Kompromiss). Dies ist ein strukturell wichtiges Ergebnis: \emph{stellare Entropieproduktion allein kann $m_e/m_p$ nicht fixieren}.

Der beobachtete Wert wird stattdessen durch mikroskopische Konsistenzeinschr\"ankungen auf atomarer und chemischer Skala eingeschr\"ankt. Chemische Bindungsenergien skalieren als $E_{\text{chem}} \sim m_e \alpha^2 c^2$; fl\"ussiges Wasser bei $T \sim 300\,$K erfordert $E_{\text{chem}} \gg k_B T$, was eine untere Grenze $m_e \gtrsim 10^{-2}\,m_{e,\text{obs}}$ ergibt. Strenger noch setzt der Bohr-Radius $a_0 = \hbar/(m_e \alpha c)$ die atomare L\"angenskala; wenn $a_0$ die intermolekularen Abst\"ande \"uberschreitet, h\"ort kondensierte Materie auf zu existieren. Diese Einschr\"ankungen erzwingen eine harte untere Grenze bei ungef\"ahr $m_e/m_p \gtrsim 10^{-5}$--$10^{-4}$.

\textbf{Zweidimensionaler Scan.} Ein gemeinsamer Scan \"uber $\alpha$ und $m_e/m_p$ stresstestet dieses mehrskalige Bild quantitativ. Das uneingeschr\"ankte 2D-Maximum liegt bei $(\alpha/\alpha_{\text{obs}}, m_e/m_{e,\text{obs}}) \approx (1{,}24;\; 0{,}30)$ mit $\mathcal{S}_{\text{total}} = 3{,}22 \times 10^{41}\,\text{J/K}$, was $\mathcal{S}(\text{obs})/\mathcal{S}_{\text{max}}^{2D} = 0{,}47$ ergibt. Das 2D-Optimum wird von der monotonen $m_e$-Abh\"angigkeit dominiert (niedrigere Elektronenmasse $\to$ geringere Opazit\"at $\to$ l\"angere Hauptreihendauer $\to$ mehr kumulative Entropie). Dies falsifiziert eine uneingeschr\"ankte einzelskalige stellare MEPP als vollst\"andiges Selektionsprinzip: die beobachtete Elektronenmasse erfordert zus\"atzliche Einschr\"ankungen aus der Atom-, Chemie- und Kernphysik. Das hierarchische Selektionsprinzip---bei dem MEPP innerhalb unabh\"angig spezifizierter Fenster operiert, die durch Konsistenzbedingungen niedrigerer Skalen definiert sind---ist daher eine notwendige Arbeitshypothese, kein durch den 2D-Scan bereits bewiesenes Ergebnis.

\textbf{Lokaler Hessian-Status.} Eine anschlie\ss{}ende 2D-Log-Hessian-Diagnostik am beobachteten Punkt, mit Koordinaten $u=\log(\alpha/\alpha_{\mathrm{obs}})$ und $v=\log(m_e/m_{e,\mathrm{obs}})$, ergibt eine steife konkave Richtung ($\lambda \approx -0{,}12344$ f\"ur $\log \mathcal{S}$) und eine nullnahe schlappe Richtung ($|\lambda|<10^{-7}$ \"uber die getesteten Schrittweiten). Dieselbe Diagnostik hat im uneingeschr\"ankten Sternmodell jedoch einen nichtverschwindenden lokalen Gradienten von ungef\"ahr $(0{,}053,\,-0{,}628)$ in $(u,v)$ am beobachteten Punkt. Der Hessian beschreibt daher die Kr\"ummung der lokalen Landschaft, zertifiziert aber \emph{keine} Stationarit\"at, KKT-Optimalit\"at oder Resolventenstabilit\"at. Daf\"ur braucht es eine unabh\"angig fixierte Constraint-Menge, einen projizierten Gradienten-Test und eine Tangential-Hessian-Pr\"ufung.

\begin{table}[htbp]
\centering
\caption{Status-Ledger f\"ur stellare Entropie- und lokale Hessian-Diagnostik. Die Tabelle trennt numerische Evidenz vom Zertifizierungsstatus, der f\"ur einen eingeschr\"ankten FST-I-Stabilit\"atsclaim n\"otig ist.}
\label{tab:kkt-ledger}
\begingroup\footnotesize
\begin{tabularx}{\textwidth}{p{0.18\textwidth}p{0.26\textwidth}p{0.18\textwidth}Y}
\toprule
\textbf{Diagnostik} & \textbf{Numerische Evidenz} & \textbf{Aktueller Status} & \textbf{N\"achstes Gate} \\
\midrule
1D-$\alpha$-Plateau & $\mathcal{S}(\alpha_{\mathrm{obs}})/\mathcal{S}_{\max}=0{,}988$; Plateau $>0{,}92$ im gangbaren Bereich & Konsistenzcheck & Selektiert $\alpha$ nicht scharf; muss mit unabh\"angigen sektor\"ubergreifenden Constraints gekoppelt werden. \\
Uneingeschr\"ankter 2D-Scan & Maximum bei $(1{,}24;\;0{,}30)$ mit $\mathcal{S}(\mathrm{obs})/\mathcal{S}_{\max}^{2D}=0{,}47$ & Negativkontrolle f\"ur einzelskalige MEPP & Beobachtetes $m_e/m_p$ muss durch vorab spezifizierte atomare und chemische Fenster begr\"undet werden. \\
Lokaler Log-Hessian & $\lambda_1\approx -0{,}12344$, $|\lambda_2|<10^{-7}$; Rohgradient $\approx(0{,}053,-0{,}628)$ & Diagnostik & Nichtverschwindender Gradient blockiert Stationarit\"at, KKT-Optimalit\"at und Resolventenstabilit\"ats-Zertifizierung. \\
Eingeschr\"ankter Tangentialtest & aktive Constraints und Projektion noch nicht fixiert & offener Zertifizierungsschritt & Constraints vorregistrieren, Gradienten projizieren und Tangential-Hessian/Resolvente testen. \\
Abgeglichene Kontrollen & No-Gamow-Variante, verschobene Constraint-Box und zuf\"allige Plateaupunkte ausstehend & noch nicht gerechnet & N\"otige Guardrail gegen posterioren Fit und modellspezifische Hessian-Artefakte. \\
\bottomrule
\end{tabularx}
\endgroup
\end{table}

Innerhalb dieses chemisch erlaubten Fensters operiert die Sternphysik nahe maximaler Entropieproduktion. Somit wird $m_e/m_p$ als durch \emph{mehrskalige Kompatibilit\"at} eingeschr\"ankt behandelt statt durch ein Einzelskalen-MEPP-Extremum: die engste Einschr\"ankung kommt von der atomaren/chemischen Skala, und gr\"o\ss{}ere Skalen (stellar, galaktisch) optimieren innerhalb dieses erlaubten Bereichs. Diese hierarchische Constraint-Lesart ist eine zentrale Arbeitshypothese des FST-Rahmens (vgl.\ FST-II \citep{FSTII}).

\textbf{Falsifizierbarkeits-Bewertung.} Unter 13 getesteten Parametervariationen ($\pm 10\%$, $\pm 20\%$, $\pm 50\%$, $+100\%$ in sowohl $\alpha$ als auch $m_e/m_p$) ergeben 7 ein $\mathcal{S} > \mathcal{S}_{\text{obs}}$ im uneingeschr\"ankten Einzelsternmodell. Die Verletzungen bei $\alpha$ sind jedoch klein (maximal 1{,}2\% \"uber dem beobachteten Wert), und alle $m_e$-Verletzungen treten bei unphysikalisch kleinen Elektronenmassen auf. Wenn strukturelle Stabilit\"atseinschr\"ankungen aus Atomphysik, Kernstabilit\"at und chemischer Gangbarkeit angewandt werden, nimmt die Anzahl der echten Verletzungen ab. Das 2D-Scan-Ergebnis ($\mathcal{S}(\text{obs})/\mathcal{S}_{\text{max}}^{2D} = 0{,}47$) demonstriert, dass das uneingeschr\"ankte MEPP-Argument unzureichend ist: die vollst\"andige FST-I-Behauptung erfordert die oben entwickelte hierarchische Einschr\"ankungsstruktur. Ein umfassender Mehrparameterscan unter Einbeziehung aller relevanten physikalischen Einschr\"ankungen ist in Vorbereitung.

\subsection{Das Higgs-VEV und kosmologische Kopplung}

\textbf{P2:} Das Higgs-VEV steht in quantitativer Beziehung zum kosmologischen S\"attigungsparameter $\Phi_0$. Beide Werte fungieren als Ordnungsparameter und Infrarot-Fixpunkte \citep{HiggsVEV2024}. Eine entdeckte Korrelation zwischen der Skala der elektroschwachen Symmetriebrechung (246 GeV) und kosmologischen Parametern w\"urde den skaleninvarianten Rahmen unterst\"utzen.

\subsection{Neutrinomassen als Optimierungsergebnis}

\textbf{P3:} Neutrinomassen sind von Null verschieden, aber minimal, weil sie als effiziente ``Entropieventile'' in Hochdichte-Kollapsprozessen fungieren.

In Kernkollaps-Supernovae m\"ussen ungef\"ahr $3 \times 10^{53}$ erg gravitativer Bindungsenergie fast ausschlie\ss{}lich durch Neutrinofluss abgestrahlt werden \citep{SupernovaNeutrinos,SupernovaSimulations}. W\"aren Neutrinos strikt masselos, w\"aren sie rein linksh\"andig (rechtsh\"andige Antineutrinos), was Flavour-Oszillationen ausschlie\ss{}en und damit die Effizienz des Energietransports zwischen verschiedenen Neutrinospezies w\"ahrend des Kollapses verringern w\"urde; w\"aren sie zu massiv, w\"urde Phasenraum-Unterdr\"uckung die Neutrinoleuchtkraft drosseln und die gesamte Energiedissipationsrate reduzieren. Die Sub-eV-Massenskala erlaubt Flavour-Oszillationen, die Energie zwischen allen drei Spezies umverteilen, die Transparenz der Neutrinosph\"are optimieren und unkontrollierte kollektive Oszillationen w\"ahrend kritischer Phasen wie des Neutronisierungsbursts verhindern \citep{NeutrinoOscillations}.

Ihre spezifischen Sub-eV-Massen sind \emph{konsistent mit} einer optimierten Konfiguration f\"ur Hochdichte-Entropieextraktion. \textbf{Konkrete Einschr\"ankung:} Die kosmologische Obergrenze f\"ur die Summe der Neutrinomassen betr\"agt $\sum m_\nu < 0{,}12\,\text{eV}$ (95\% C.L., Planck 2018 \citep{Planck2018}). Wir beobachten, dass der FST-I-Rahmen mit Massen nahe der unteren Grenze des entropieeffizienten Fensters konsistent ist; dies ist jedoch eine qualitative Post-hoc-Beobachtung und keine quantitative Vorhersage. Eine pr\"azise Vorhersage der Massenhierarchie (normal vs.\ invertiert) oder der leichtesten Neutrinomasse erfordert eine quantitative Berechnung der K\"uhlungsoptimierung, die noch nicht verf\"ugbar ist. Dies ist testbar mit JUNO, DUNE und zuk\"unftigen kosmologischen Durchmusterungen. In demselben heuristischen K\"uhlventil-Bild w\"aren stark gekoppelte sterile Neutrinos unplausibel, falls sie die optimierte K\"uhlung robust st\"orten; dies ist aber noch keine unabh\"angige FST-I-Vorhersage.

\subsection{Ehrliche Abgrenzung}

Der Rahmen sagt derzeit nicht vorher:
\begin{itemize}
    \item Die exakte Anzahl der Fermionengenerationen
    \item Die detaillierte Flavour-Mischungsstruktur (CKM- und PMNS-Matrizen)
    \item Die spezifische mikroskopische Natur von Dunkle-Materie-Kandidaten
\end{itemize}

Dies spiegelt wider, dass MEPP den gangbaren Parameterraum einschr\"ankt, aber nicht notwendigerweise jeden Freiheitsgrad eindeutig diktiert, wenn mehrere Konfigurationen entartete Payoffs ergeben.

\subsection{Strenge Falsifizierbarkeit}

\begin{quote}
\textit{Wenn eine theoretische Variation der Standardmodell-Parameter entdeckt wird, die das integrierte Entropieproduktionsfunktional $\mathcal{S}$ signifikant erh\"oht, w\"ahrend gleichzeitig alle strukturellen Stabilit\"atseinschr\"ankungen (nukleare Stabilit\"at, atomare Bindung, stellare Langlebigkeit und chemische Komplexit\"at) erhalten bleiben, dann ist die Hypothese, dass das Standardmodell eine stabilit\"atsselektierte thermodynamische Konfiguration darstellt, falsifiziert. Das semi-analytische Sternmodell (Abschnitt~\ref{sec:stellar_model}) zeigt, dass die uneingeschr\"ankte Entropieproduktion bei $\alpha/\alpha_{\mathrm{obs}} \approx 1{,}6$ ihren H\"ochstwert erreicht, wobei nur marginal h\"ohere Werte erzielt werden. Die beobachtete Feinstrukturkonstante liegt innerhalb des breiten Plateaus ($\mathcal{S}/\mathcal{S}_{\max} > 0{,}96$ f\"ur $\alpha/\alpha_{\mathrm{obs}} \in [0{,}75;\; 3{,}0]$), konsistent mit Stabilit\"atsselektion aus einer entarteten Mannigfaltigkeit statt scharfer Optimierung.}
\end{quote}

\noindent Konkret w\"are eine falsifizierende Variation eine Parameterst\"orung $\delta p$ im Standardmodell-Parameterraum, die (i)~alle bekannten Stabilit\"atseinschr\"ankungen erhalt (Vakuumstabilit\"at, BBN-Kompatibilit\"at, stellare Z\"undung), (ii)~die integrierte Entropieproduktion erh\"oht, $\mathcal{S}[p + \delta p] > \mathcal{S}[p]$, und (iii)~nicht durch bestehende experimentelle Schranken ausgeschlossen ist. Die Identifikation einer solchen St\"orung---oder ein Beweis, dass keine existiert---stellt die zentrale offene Herausforderung dieses Programms dar.

\textbf{Konkretes quantitatives Kriterium.} Um zu verhindern, dass die strukturellen Stabilit\"atseinschr\"ankungen als unfalsifizierbare Ausweichklausel dienen, legen wir uns auf den folgenden vorregistrierten Test fest: Wenn der Mehrparameter-Scan (Schlussfolgerung, Punkt~1) \"uber die f\"unf Kernparameter $\{\alpha_{EM}, \alpha_s, G_F, m_e/m_p, m_u/m_d\}$ abgeschlossen ist und mehr als zwei dieser f\"unf Parameter au\ss{}erhalb des 80\%-Plateaus von $\mathcal{S}/\mathcal{S}_{\max}$ liegen (d.\,h.\ $\mathcal{S}(\text{obs})/\mathcal{S}_{\max} < 0{,}80$ entlang des jeweiligen Einparameter-Querschnitts, w\"ahrend alle anderen Parameter auf ihren beobachteten Werten gehalten werden und alle strukturellen Stabilit\"atseinschr\"ankungen erf\"ullt sind), dann ist die Entropiedominanz-Konjektur in ihrer aktuellen Form falsifiziert. Die Liste der strukturellen Stabilit\"atseinschr\"ankungen ist auf die aus etablierter Kern- und Atomphysik aufz\"ahlbaren Constraints festgelegt: Vakuumstabilit\"at, BBN-Ertrag innerhalb von $2\sigma$ der beobachteten H\"aufigkeiten, stellare Z\"undung mit $T_c > 4 \times 10^6\,\text{K}$ und chemische Bindungsenergien $E_{\text{chem}} \gg k_B T_{\text{surface}}$. Eine nachtr\"agliche Erg\"anzung neuer Constraints ist nicht zul\"assig.

\subsection{Einheitliches Stabilit\"atsprotokoll}\label{sec:stability-protocol}

Die obigen Falsifizierbarkeitskriterien erfordern ein konkretes, reproduzierbares
Testverfahren. Das folgende Protokoll operationalisiert das
Resolventenstabilit\"atskonzept f\"ur jedes System innerhalb des
FST-Rahmens, mit dom\"anenspezifischen Instanziierungen in FST-I
(Teilchenphysik), FST-II (chemische Netzwerke) und FST-III
(Proteinfaltung).

\begin{definition}[FST-Stabilit\"atstest]\label{def:stability-test}
Gegeben ein System, parametrisiert durch $p \in \mathcal{P}$, besteht der
\emph{FST-Stabilit\"atstest} aus drei Schritten:
\begin{enumerate}
  \item \textbf{St\"orung.} Anwendung einer St\"orung $\delta p$ auf die
        Gleichgewichtskonfiguration $p^\star$. Die Natur von $\delta p$
        ist dom\"anenspezifisch:
        \begin{itemize}
          \item FST-I: $\delta p$ repr\"asentiert Variationen der
                Standardmodell-Parameter (Kopplungskonstanten, Massen,
                Mischungswinkel).
          \item FST-II: $\delta p$ repr\"asentiert \"Anderungen der
                Netzwerktopologie oder Reaktionsratenkonstanten.
          \item FST-III: $\delta p$ repr\"asentiert Konformations\"anderungen
                 oder Punktmutationen in der Aminos\"auresequenz.
         \end{itemize}
  \item \textbf{Prim\"are Diagnostik.} Berechnung des Spektralradius
         $\rho(J)$ des Best-Response-Jacobians
         $J = I - \eta\, H$ bei $p^\star$, wobei $\eta > 0$ ein
         Lernratenparameter ist, der die Schrittweite der
         Best-Response-Dynamik kontrolliert (f\"ur hinreichend kleines
         $\eta$ ist das Spektralradius-Kriterium $\rho(J)<1$ \"aquivalent
         zur Negativdefinitheit von~$H$; siehe FST-III f\"ur eine
         empirische Demonstration, dass das Urteil \"uber einen weiten
         Bereich $\eta$-insensitiv ist)\footnote{Vorzeichenkonvention:
         Hier bezeichnet $H = -\nabla^2 S$ den \emph{negativen} Hessian
         des Potentials, sodass $H$ an einem lokalen Maximum von~$S$
         positiv definit ist. Mit dieser Konvention hat
         $J = I - \eta H$ die Eigenwerte $1 - \eta \lambda_i$ mit
         $\lambda_i > 0$, woraus $\rho(J) < 1$ f\"ur
         $0 < \eta < 2/\lambda_{\max}$ folgt. Das Stabilit\"atsurteil
         (Positivdefinitheit von~$H$) ist von~$\eta$ unabh\"angig; nur
         die Konvergenzrate der diskreten Iteration h\"angt von der
         Schrittweite ab.} und $H$ der aus dem Potential abgeleitete
         Stabilit\"atsoperator ist (Annahme~\ref{ass:potential-game}).
         Eichrichtungen werden durch Projektion auf den physikalischen
         Unterraum ausgeschlossen.
  \item \textbf{Sekund\"are Diagnostik.} Auswertung der
         Entropieproduktionsrate $\langle\dot{S}\rangle$ als
         \emph{Korrelat}, nicht als Definition, von Stabilit\"at.
\end{enumerate}
F\"ur den vorregistrierten Kleinschrittbereich von $\eta$ ist das System nur
dann \emph{resolvent-stabil}, wenn $\rho(J) < 1$ f\"ur alle Nicht-Eich-Moden
gilt. \"Aquivalent muss der projizierte Operator $H$ auf dem physikalischen
Tangentialraum positiv sein, und die gew\"ahlten Schrittweiten m\"ussen
$0<\eta\lambda_i<2$ f\"ur die entsprechenden Nicht-Eich-Eigenwerte erf\"ullen.
\end{definition}

\begin{remark}[Guardrail f\"ur eingeschr\"ankte Optima]\label{rem:kkt-guardrail}
Bei eingeschr\"ankten Parameterproblemen wie FST-I ist Hessian-Definitheit
allein kein Stabilit\"atszertifikat. Die aktiven Constraints m\"ussen vor dem
Test fixiert sein, der Gradient am beobachteten Punkt muss entweder im
physikalischen Tangentialraum verschwinden oder die entsprechenden
Karush--Kuhn--Tucker-Bedingungen erf\"ullen, und der Hessian muss danach auf
dem Tangentialraum statt auf nachtr\"aglich ausgeschlossenen Richtungen
ausgewertet werden. Eine Kr\"ummungsmatrix an einem nicht-station\"aren Punkt
ist nur eine Landschaftsdiagnostik. Sie kann ein Stabilit\"atsprogramm
motivieren, etabliert aber keine Resolventenstabilit\"at.
\end{remark}

\begin{remark}[Falsifizierbarkeit]\label{rem:falsifiability-protocol}
Der FST-Stabilit\"atstest liefert zwei testbare Vorhersagen:
\begin{enumerate}[label=(\roman*)]
  \item Beobachtete Gleichgewichtskonfigurationen sind resolvent-stabil,
        d.\,h.\ $\rho(J) < 1$ in allen Nicht-Eich-Richtungen.
  \item Stabilit\"at und Entropieproduktion sind \"uber den zug\"anglichen
        Parameterraum positiv korreliert.
\end{enumerate}
Ein systematisches Fehlen dieser Korrelation oder ein persistentes
$\rho(J) > 1$ an einer beobachteten Gleichgewichtskonfiguration w\"urde
den Rahmen falsifizieren. F\"ur eine explizite Demonstration im
chemischen Kontext siehe das Drei-Spezies-Autokatalyse-Modell in
FST-II~\cite{FSTII}.
\end{remark}

\subsection{Rechnerische Verifikation: Feinstrukturkonstante}
\label{sec:comp-verification-alpha}

Als unabh\"angige Pr\"ufung des semi-analytischen Sternmodells (Abschnitt~\ref{sec:stellar_model}) implementierten wir ein Gamow-erweitertes numerisches Modell, das die integrierte Entropieproduktion $\mathcal{S}_{\text{total}}$ als stetige Funktion von $\alpha$ auswertet und das Maximum mittels gradientenbasierter Optimierung lokalisiert. Zwei Modellvarianten wurden verglichen: ein \emph{einfaches Skalierungsmodell}, das die Leuchtkraft als Potenzgesetz $L \propto \alpha^4$ ohne kernphysikalische Korrekturen behandelt, und ein \emph{Gamow-erweitertes Modell}, das den $\alpha$-abh\"angigen Gamow-Tunnelfaktor f\"ur pp-Ketten-Fusion einbezieht. Die Ergebnisse sind in Tabelle~\ref{tab:comp-verification-alpha} zusammengefasst.

\begin{table}[ht]
\centering
\caption{Rechnerische Verifikation des Entropieproduktions-Maximums bez\"uglich der Feinstrukturkonstante. Die dimensionslose Koordinate $x_\alpha = \alpha / \alpha_{\text{obs}}$ parametrisiert den Scan; $x_{\max}$ bezeichnet die Lage des Entropiemaximums; $\mathcal{S}_{\text{obs}} / \mathcal{S}_{\max}$ ist das Verh\"altnis der beobachteten Entropieproduktion zu ihrem Maximalwert; $H(\alpha_{\text{obs}})$ ist die Hesse-Matrix (zweite Ableitung) am beobachteten Wert. Skript: \texttt{stellar\_testbed.py}.}
\label{tab:comp-verification-alpha}
\begin{tabular}{lccc}
\toprule
\textbf{Modell} & $x_{\max}$ & $\mathcal{S}_{\text{obs}} / \mathcal{S}_{\max}$ & $H(\alpha_{\text{obs}})$ \\
\midrule
Einfache Skalierung & 0{,}658 & 43{,}3\% & $+6{,}75$ (Minimum) \\
Gamow-erweitert & 1{,}605 & 98{,}75\% & $-0{,}18$ (Maximum) \\
\bottomrule
\end{tabular}
\end{table}

Das Gamow-erweiterte Modell platziert das Entropieproduktions-Maximum bei $x_\alpha = 1{,}605$, in \"Ubereinstimmung mit dem Ergebnis der semi-analytischen Tabelle~\ref{tab:alpha_scan} ($x_\alpha \approx 1{,}5$--$1{,}6$). Die beobachtete Feinstrukturkonstante erreicht 98{,}75\% dieses Maximums. In diesem Pilotmodell hat die Hesse-Matrix bei $\alpha_{\text{obs}}$ den Wert $H = -0{,}18$, was zeigt, dass der beobachtete Wert auf der konkaven (Maximum-) Seite der Entropielandschaft liegt. Das einfache Skalierungsmodell, das die Gamow-Tunnelabh\"angigkeit von $\alpha$ wegl\"asst, ergibt ein qualitativ anderes Bild: das Maximum verschiebt sich zu $x_\alpha = 0{,}658$ und der beobachtete Wert erreicht nur 43{,}3\% von $\mathcal{S}_{\max}$, mit positiver Hesse-Matrix, die ein lokales Minimum anzeigt. Dieser Vergleich zeigt, dass die Kernphysik---konkret die exponentielle Empfindlichkeit der Tunnelrate auf $\alpha$---wesentlich f\"ur das MEPP-Argument ist: der Gamow-Faktor erzeugt das breite Maximum nahe $\alpha_{\text{obs}}$, das die Entropiedominanz-Hypothese nicht-trivial macht.

Das Gamow-erweiterte Modell liefert damit einen ersten quantitativen Konsistenztest: die beobachtete Feinstrukturkonstante liegt innerhalb von 1{,}25\% des Entropieproduktions-Maximums ($\mathcal{S}_{\text{obs}}/\mathcal{S}_{\max} = 0{,}9875$), mit negativer Kr\"ummung $H = -0{,}18$, die zeigt, dass der beobachtete Wert auf der konkaven Schulter des breiten Maximums liegt. Da das Maximum selbst bei $x_\alpha=1{,}605$ liegt, ist diese Kr\"ummung kein Stationarit\"atszertifikat bei $\alpha_{\text{obs}}$. Obwohl dieses Ergebnis modellabh\"angig ist (es beruht auf dem polytropischen Sternmodell und der vereinfachten Gamow-Behandlung), ist die qualitative Schlussfolgerung---dass $\alpha_{\text{obs}}$ auf einem breiten Hochentropie-Plateau liegt, wenn nukleares Tunneln ber\"ucksichtigt wird---eine robuste Diagnostik innerhalb dieses Pilotmodells.

% ===================================================================
\section{Kosmologische Implikationen und der Zeitpfeil}
\label{sec:cosmology}
% ===================================================================

Das thermodynamische Spiel, das die fundamentalen Teilchensektoren bestimmt, entfaltet sich vor der dynamisch expandierenden FLRW-Metrik. Wir vermerken---ohne dabei Erkl\"arungskraft zu beanspruchen---dass sp\"ate kosmische Beschleunigung \emph{qualitativ konsistent} mit Entropiemaximierung ist: Die Ausweitung des Phasenraumvolumens erh\"alt die Kapazit\"at f\"ur fortgesetzte Dissipation, wenn lokalisierter nuklearer Brennstoff ersch\"opft ist. Diese Beobachtung ist nicht neu (vgl.\ \citep{EntropyMaximum} f\"ur eine quantitative Diskussion kosmologischer Entropiebudgets) und stellt keine Vorhersage der kosmologischen Konstante aus FST-I-Prinzipien dar. Eine quantitative Herleitung wurde nicht erreicht; dies bleibt spekulativ (Tabelle~\ref{tab:status-claims}).

\subsection{Stabilit\"at vs.\ Maximierung}

Die flache Entropielandschaft (Tabelle~\ref{tab:alpha_scan}: $\mathcal{S}/\mathcal{S}_{\max} > 0{,}92$ \"uber den gangbaren Bereich) legt nahe, dass das Fine-Tuning-Problem besser als Stabilit\"atsproblem denn als Maximierungsr\"atsel formuliert werden sollte. Viele Konfigurationen liefern vergleichbare Entropieproduktion erster Ordnung; die beobachtete Konfiguration k\"onnte durch Kreuzkopplungen zweiter Ordnung zwischen Parametersektoren selektiert werden. Diese Hypothese der ``Stabilit\"atsselektion'' bietet eine Alternative sowohl zum Landscape-Ansatz als auch zum anthropischen Argument, bleibt jedoch eine \emph{strukturelle Konjektur}. Die neue 2D-Log-Hessian-Diagnostik zeigt eine steife konkave Richtung und eine nullnahe Richtung, w\"ahrend der uneingeschr\"ankte Gradient nicht verschwindet. Ob der eingeschr\"ankte Tangential-Hessian in allen unabh\"angig zul\"assigen Nicht-Eich-Richtungen negativ definit ist---zusammen mit den erforderlichen KKT-Bedingungen---ist eine offene Frage, deren Kl\"arung den in der Schlussfolgerung als Priorit\"at~2 aufgef\"uhrten Mehrparameter-Scan erfordert.

\subsection{Status der Behauptungen}

Tabelle~\ref{tab:status-claims} klassifiziert die zentralen Behauptungen dieses Beitrags nach epistemischem Status.

\begin{table}[htbp]
\centering
\caption{Status der Behauptungen. \textsc{Etabliert}: bewiesene oder weithin akzeptierte Ergebnisse, die hier zitiert werden; \textsc{Konjektur}: neue, aber testbare Behauptungen; \textsc{Spekulativ}: Rahmenvorschl\"age ohne direkten empirischen Test.}
\label{tab:status-claims}
\small
\begin{tabularx}{\textwidth}{@{}>{\raggedright\arraybackslash}p{4.0cm}>{\raggedright\arraybackslash}p{2.7cm}Y@{}}
\toprule
\textbf{Behauptung} & \textbf{Status} & \textbf{Evidenz / Referenz} \\
\midrule
Nash-Gleichgewicht / Mean-Field-Spieltheorie & \textsc{Etabliert} & Nash (1950); Lasry \& Lions (2007) \\
MEPP als heuristisches Organisationsprinzip & \textsc{Etabliert} & Martyushev \& Seleznev (2006); Dyke \& Kleidon (2010) \\
Born-Regel-Neuinterpretation \"uber Envariance & \textsc{Konjektur} & Born-Regel ist Standard; FST-I nutzt Zureks Envariance-Route als umstrittene konzeptuelle Analogie \\
Diproton-Katastrophe \& Triple-Alpha-Sensitivit\"at & \textsc{Etabliert} & Adams (2019) \\
Quanten--Strategie-Korrespondenz (kommutative Spezialisierung) & \textsc{Konjektur} & Abschn.~\ref{sec:quantum}; diagonale Reduktion von Lindblad-OT auf Lasry--Lions; Koh\"arenzen erfordern nichtkommutative Erweiterung \\
SM-Parameter als Nash-Gleichgewicht des Entropiespiels & \textsc{Konjektur} & Conjecture~\ref{conj:entropy-dominance}; nur Pilotstudie \\
Stellare Entropie peakt nahe $\alpha_{\text{obs}}$ (98,8\%) & \textsc{Konjektur} & Tabelle~\ref{tab:alpha_scan}; Einkanal-, Einstern-Modell \\
Resolventenstabilit\"at als Fine-Tuning-Erkl\"arung & \textsc{Konjektur} & Abschn.~\ref{sec:cosmology}; strukturelle Analogie, kein eingeschr\"anktes Tangential-Hessian-/KKT-Zertifikat berechnet \\
2D-Log-Hessian-Pilot $\alpha$/$m_e$ & \textsc{Konjektur} & Eine steife konkave Richtung und eine nullnahe Richtung; nichtverschwindender Gradient, kein KKT-Zertifikat, keine volle Resolventenstabilit\"at \\
Symmetriebrechung als Phasen\"ubergang zwischen Gleichgewichten & \textsc{Konjektur} & Abschn.~\ref{sec:symmetry}; Neuinterpretation bekannter Physik \\
Neutrinomassen als Entropieventil-Optimierung & \textsc{Spekulativ} & P3; keine quantitative Massenvorhersage \\
Higgs-VEV--kosmologische S\"attigungs-Verbindung & \textsc{Spekulativ} & P2; keine Ableitung von $\Lambda$ aus MEPP \\
Kosmologische Beschleunigung als Entropiestrategie & \textsc{Spekulativ} & Abschn.~\ref{sec:cosmology}; nur heuristisch \\
\bottomrule
\end{tabularx}
\end{table}

% ===================================================================
\section{Schlussfolgerung}
\label{sec:conclusion}
% ===================================================================

Diese Arbeit schl\"agt einen vereinigenden \emph{interpretativen Rahmen} f\"ur die fundamentale Physik vor, der auf einem einzigen organisierenden Prinzip basiert: thermodynamische Optimierung unter strategischer Interaktion. Innerhalb dieses Rahmens wird das Standardmodell nicht als willk\"urliche Sammlung von Parametern betrachtet, sondern als m\"oglicherweise konsistent mit dem eingeschr\"ankten Gleichgewichtsergebnis eines mehrskaligen Optimierungsspiels. Wir betonen, dass FST-I in diesem Stadium eine strukturierte Hypothese und keine pr\"adiktive Theorie ist: Der Ansatz erzeugt keine quantitativen Vorhersagen jenseits etablierter Physik und leitet keine Parameterwerte aus ersten Prinzipien ab. Sein Beitrag ist ein konzeptuelles Vokabular, das bekannte Parametersensitivit\"aten in einer spieltheoretischen Architektur organisiert.

Auf Teilchenebene wird die Quantenmechanik als Gleichgewicht gemischter Strategien uminterpretiert, wobei das Pfadintegral als Nash-Selektor-Analogie fungiert, die nicht-optimale Strategien durch destruktive Interferenz unterdr\"uckt. Die Born-Regel wird \"uber Envariance als m\"ogliche Gleichgewichtsverteilung umgedeutet---eine konzeptuelle Uminterpretation, keine neue Herleitung. Symmetriebrechung wird als Phasen\"ubergang zwischen Gleichgewichten uminterpretiert, bei dem ver\"anderte \"au\ss{}ere Bedingungen symmetrische Konfigurationen destabilisieren und das System in Richtung entropieproduzierender Attraktoren treiben. Fine-Tuning wird in dieser Sicht als m\"ogliche Signatur eingeschr\"ankter Optimierung statt als blo\ss{}er Zufall behandelt.

Entscheidend ist, dass dieser Rahmen die formale Struktur der etablierten Physik nicht modifiziert. Es werden keine neuen Teilchen, Kr\"afte oder dynamischen Gesetze eingef\"uhrt. Stattdessen liegt der Beitrag in der Identifizierung einer gemeinsamen Optimierungslogik, die bekannte Ph\"anomene \"uber Skalen hinweg organisiert. Die Standardmodell-Parameter werden als m\"oglicherweise stabilit\"atskompatibles Werkzeugset f\"ur Entropieproduktion in einem Universum interpretiert, das durch thermodynamische Einschr\"ankungen regiert wird.

Der Ansatz umfasst explizite Falsifizierbarkeitskriterien (Abschnitt~\ref{sec:predictions}), einschlie\ss{}lich eines vorregistrierten quantitativen Kriteriums f\"ur den geplanten Mehrparameter-Scan. Eine erste semi-analytische Pilotstudie zeigt, dass die beobachtete Feinstrukturkonstante innerhalb eines breiten Plateaus nahezu maximaler stellarer Entropieproduktion liegt ($\mathcal{S}/\mathcal{S}_{\max} = 0{,}988$, aber mit $\mathcal{S}/\mathcal{S}_{\max} > 0{,}92$ \"uber den gesamten lebensf\"ahigen Bereich)---ein Konsistenzcheck, kein entscheidender Test. Die Flachheit dieses Plateaus bedeutet, dass das 1D-Ergebnis zwar f\"ur die Tragf\"ahigkeit des Rahmens notwendig ist, aber keine scharfe Parameterselektion demonstriert. Ein zweidimensionaler Scan ergibt $\mathcal{S}/\mathcal{S}_{\max}^{2D} = 0{,}47$ und zeigt damit, dass eine einzelskalige stellare MEPP-Argumentation als eigenst\"andiges Selektionsprinzip scheitert: Ohne hierarchische mehrskalige Constraints aus Atom- und chemischer Physik scheitert die Entropiedominanz-Vermutung bereits in zwei Dimensionen. Die volle mehrskalige Constraint-Hierarchie ist daher eine erforderliche Arbeitshypothese von FST-I, ihre unabh\"angige KKT-/Resolventenvalidierung bleibt aber offen. Diese Pilotstudie auf den gesamten Standardmodell-Parameterraum zu erweitern, bleibt die zentrale rechnerische Herausforderung.

Dieser Beitrag repr\"asentiert den ersten Schritt eines umfassenderen Forschungsprogramms. Die unmittelbar n\"achsten Schritte sind, nach Priorit\"at geordnet:
\begin{enumerate}
    \item \textbf{Mehrparameter-Stellarscan (kritisch):} Erweiterung des semi-analytischen Modells auf gleichzeitige Variation von $\alpha_s$, $G_F$ und $m_u/m_d$ unter Einbeziehung der Kernstabilit\"atsgrenzen aus \citep{Adams2019}. Dies ist der direkteste Weg, die Entropiedominanz-Vermutung \"uber die hier pr\"asentierten Einzelparameter-Ergebnisse hinaus zu testen.
    \item \textbf{KKT- und Hesse-Analyse (kritisch):} Berechnung des projizierten Gradienten, der aktiven Constraints und der Stabilit\"atsmatrix zweiter Ordnung $\partial^2 \mathcal{S}/\partial p_i \partial p_j$ im Tangentialraum am beobachteten Parameterpunkt, um die Resolventenstabilit\"ats-Vermutung quantitativ zu testen. Ein negativ-definiter Tangential-Hessian nach unabh\"angig spezifizierter KKT-Pr\"ufung w\"are der st\"arkste Nachweis f\"ur den Rahmen.
    \item \textbf{Mehrkanalentropie:} Einbeziehung gravitativer (Strukturbildung), chemischer (molekulare Komplexit\"at) und nuklearer (Nukleosynthese-Ausbeute) Entropiekan\"ale \"uber die Sternstrahlung hinaus.
    \item \textbf{Vollst\"andige Sternentwicklungscodes:} Ersetzung des polytropischen Modells durch MESA-artige Sternentwicklungssimulationen mit nicht-standardm\"a\ss{}igen nuklearen Parametern.
\end{enumerate}

Wenn das oben skizzierte Berechnungsprogramm erfolgreich demonstriert, dass die beobachteten Parameter ein mehrdimensionales Stabilit\"atsmaximum besetzen, w\"urde dieses Bild nahelegen, dass physikalische Gesetze auf verschiedenen Skalen eine gemeinsame Optimierungslogik teilen, anstatt unabh\"angig voneinander konstruiert zu sein. In diesem Fall w\"urde Fine-Tuning weniger als zu erkl\"arendes Mysterium erscheinen, sondern als Diagnose: eine Spur des Gleichgewichts in einem thermodynamischen Spiel, das vom Quantenvakuum bis zu den gr\"o\ss{}ten kosmischen Strukturen reicht. Bis dahin bleibt FST-I ein vielversprechender, aber unvalidierter interpretativer Rahmen.

% ===================================================================
% Literaturverzeichnis
% ===================================================================

\bibliographystyle{plainnat}
{\small
\setlength{\parskip}{0pt}
\begin{thebibliography}{99}
\setlength{\itemsep}{0pt}

\bibitem[Aghanim et al.(2020)]{Planck2018}
Aghanim, N. et al. (Planck Collaboration) (2020).
\newblock Planck 2018 results. VI. Cosmological parameters.
\newblock \textit{Astronomy \& Astrophysics}, 641, A6.
\newblock arXiv:1807.06209.

\bibitem[Adams(2019)]{Adams2019}
Adams, F.~C. (2019).
\newblock The degree of fine-tuning in our universe---and others.
\newblock \textit{Physics Reports}, 807, 1--111.
\newblock arXiv:1902.03928.

\bibitem[Crooks(1999)]{Crooks1999}
Crooks, G.~E. (1999).
\newblock Entropy production fluctuation theorem and the nonequilibrium work relation for free energy differences.
\newblock \textit{Physical Review E}, 60, 2721--2726.

\bibitem[England(2013)]{England2013}
England, J.~L. (2013).
\newblock Statistical physics of self-replication.
\newblock \textit{The Journal of Chemical Physics}, 139, 121923.
\newblock DOI: 10.1063/1.4818538.

\bibitem[Dewar(2003)]{Dewar2003}
Dewar, R.~C. (2003).
\newblock Information theory explanation of the fluctuation theorem, maximum entropy production and self-organized criticality in non-equilibrium stationary states.
\newblock \textit{Journal of Physics A: Mathematical and General}, 36(3), 631.

\bibitem[Djete and Touzi(2026)]{DjeteTouzi2026}
M.~F.~Djete and N.~Touzi (2026).
\newblock On Approximate Nash Equilibria in Mean Field Games.
\newblock arXiv:2601.20910.

\bibitem[Grinstein \& Linsker(2007)]{GrinsteinLinsker2007}
Grinstein, G. \& Linsker, R. (2007).
\newblock Comments on a derivation and application of the ``maximum entropy production'' principle.
\newblock \textit{Journal of Physics A: Mathematical and Theoretical}, 40(31), 9717.

\bibitem[Bousso \& Polchinski(2000)]{Bousso2000}
Bousso, R. \& Polchinski, J. (2000).
\newblock Quantization of four-form fluxes and dynamical neutralization of the cosmological constant.
\newblock \textit{JHEP}, 06, 006.

\bibitem[Valentini \& Westman(2005)]{ValentiniWestman2005}
Valentini, A. \& Westman, H. (2005).
\newblock Dynamical origin of quantum probabilities.
\newblock \textit{Proc.\ R.\ Soc.\ A}, 461(2053), 253--272. arXiv:quant-ph/0403034.

\bibitem[Friston(2010)]{Friston2010}
Friston, K. (2010).
\newblock The free-energy principle: a unified brain theory?
\newblock \textit{Nature Reviews Neuroscience}, 11, 127--138.

\bibitem[Friston et al.(2006)]{FristonFreeEnergy}
Friston, K., Kilner, J. \& Harrison, L. (2006).
\newblock A free energy principle for the brain.
\newblock \textit{Journal of Physiology -- Paris}, 100, 70--87.

\bibitem[Kleidon \& Lorenz(2004)]{MEPP2004}
Kleidon, A. \& Lorenz, R.~D. (eds.) (2004).
\newblock Non-equilibrium Thermodynamics and the Production of Entropy: Life, Earth, and Beyond.
\newblock \textit{Springer}.

\bibitem[Dyke \& Kleidon(2010)]{MEPPFoundations}
Dyke, J. \& Kleidon, A. (2010).
\newblock The Maximum Entropy Production Principle: Its Theoretical Foundations and Applications to the Earth System.
\newblock \textit{Entropy}, 12, 613--630.
\newblock DOI: 10.3390/e12030613.

\bibitem[Abel(2023)]{Abel2023}
Abel, C.~R. (2023).
\newblock The quantum foundations of utility and value.
\newblock \textit{Phil.\ Trans.\ R.\ Soc.\ A}, 381, 20220286.

% Maity (2024) removed from active citations; retained in working bibliography for reference.
% Original: arXiv:2402.13395 (not peer-reviewed, non-physics venue).

\bibitem[Lindblad(1976)]{Lindblad1976}
Lindblad, G. (1976).
\newblock On the generators of quantum dynamical semigroups.
\newblock \textit{Communications in Mathematical Physics}, 48, 119--130.

\bibitem[Martyushev \& Seleznev(2006)]{MartyushevSeleznev2006}
Martyushev, L.~M. \& Seleznev, V.~D. (2006).
\newblock Maximum entropy production principle in physics, chemistry and biology.
\newblock \textit{Physics Reports}, 426, 1--45.

\bibitem[Martyushev(2010)]{Martyushev2010}
Martyushev, L.~M. (2010).
\newblock The maximum entropy production principle: two basic questions.
\newblock \textit{Phil.\ Trans.\ R.\ Soc.\ B}, 365(1545), 1333--1334.

\bibitem[Lasry and Lions(2006)]{LasryLions2006}
J.-M.~Lasry and P.-L.~Lions (2006).
\newblock Jeux \`a champ moyen.\ I -- Le cas stationnaire.
\newblock \textit{Comptes Rendus Math\'ematique}, 343(9), 619--625.
\newblock DOI: 10.1016/j.crma.2006.09.019.

\bibitem[Lasry \& Lions(2007)]{LasryLions2007}
Lasry, J.-M. \& Lions, P.-L. (2007).
\newblock Mean field games.
\newblock \textit{Japanese Journal of Mathematics}, 2(1), 229--260.

\bibitem[Padmanabhan(2010)]{Padmanabhan2010}
Padmanabhan, T. (2010).
\newblock Thermodynamical aspects of gravity: New insights.
\newblock \textit{Rep. Prog. Phys.}, 73, 046901.
\newblock DOI: 10.1088/0034-4885/73/4/046901.

\bibitem[Feynman \& Hibbs(1965)]{PathIntegral}
Feynman, R.~P. \& Hibbs, A.~R. (1965).
\newblock \textit{Quantum Mechanics and Path Integrals}.
\newblock McGraw-Hill.

\bibitem[Kleinert(2009)]{PathIntegralOptimization}
Kleinert, H. (2009).
\newblock \textit{Path Integrals in Quantum Mechanics, Statistics, Polymer Physics, and Financial Markets} (5th ed.).
\newblock World Scientific.

% Pramanik (2024) removed from active citations; retained in working bibliography for reference.
% Original: Ada Academica (non-standard venue, not indexed in Scopus/WoS).

\bibitem[Busemeyer \& Bruza(2012)]{BusemeyerBruza2012}
Busemeyer, J.~R. \& Bruza, P.~D. (2012).
\newblock \textit{Quantum Models of Cognition and Decision}.
\newblock Cambridge University Press.

\bibitem[Eisert et al.(1999)]{QuantumOptimization}
Eisert, J., Wilkens, M. \& Lewenstein, M. (1999).
\newblock Quantum games and quantum strategies.
\newblock \textit{Physical Review Letters}, 83(15), 3077--3080.

\bibitem[Chakrabarti et al.(2024)]{HiggsVEV2024}
Chakrabarti, S., Anagha V., Ganesh, S. \& Menon, V. (2024).
\newblock A cosmological reconstruction of the {H}iggs vacuum expectation value.
\newblock \textit{Eur.\ Phys.\ J.\ C}, 84, 1250.
\newblock DOI: 10.1140/epjc/s10052-024-13626-4. arXiv:2409.15962.

\bibitem[Schlosshauer(2007)]{MeasurementClassical}
Schlosshauer, M. (2007).
\newblock \textit{Decoherence and the Quantum-to-Classical Transition}.
\newblock Springer.

\bibitem[Jaynes(1957)]{Jaynes1957}
Jaynes, E.~T. (1957).
\newblock Information theory and statistical mechanics.
\newblock \textit{Physical Review}, 106(4), 620--630.

\bibitem[Janka(2012)]{SupernovaNeutrinos}
Janka, H.-T. (2012).
\newblock Explosion mechanisms of core-collapse supernovae.
\newblock \textit{Annual Review of Nuclear and Particle Science}, 62, 407--451.

\bibitem[Burrows(2013)]{SupernovaSimulations}
Burrows, A. (2013).
\newblock Colloquium: Perspectives on core-collapse supernova theory.
\newblock \textit{Reviews of Modern Physics}, 85(1), 245--261.

\bibitem[Johns et al.(2025)]{NeutrinoOscillations}
Johns, L., Richers, S. \& Wu, M.-R. (2025).
\newblock Neutrino oscillations in core-collapse supernovae and neutron star mergers.
\newblock \textit{Ann.\ Rev.\ Nucl.\ Part.\ Sci.}, 75, 399--423.
\newblock DOI: 10.1146/annurev-nucl-121423-100853. arXiv:2503.05959.

\bibitem[Smolin(1992)]{Smolin1992}
Smolin, L. (1992).
\newblock Did the Universe evolve?
\newblock \textit{Classical and Quantum Gravity}, 9, 173--191.

\bibitem[Susskind(2005)]{Susskind2005}
Susskind, L. (2005).
\newblock \textit{The Cosmic Landscape}.
\newblock Little, Brown.

\bibitem[Tegmark et al.(2006)]{Tegmark2006}
Tegmark, M. et al. (2006).
\newblock Dimensionless constants, cosmology, and other dark matters.
\newblock \textit{Phys. Rev. D}, 73, 023505.

\bibitem['t~Hooft(2016)]{tHooft2016}
't~Hooft, G. (2016).
\newblock \textit{The Cellular Automaton Interpretation of Quantum Mechanics}.
\newblock Springer, Fundamental Theories of Physics, Vol.~185.

\bibitem[Busch et al.(2007)]{HeisenbergCovariant}
Busch, P., Heinonen, T. \& Lahti, P. (2007).
\newblock Heisenberg's uncertainty principle.
\newblock \textit{Physics Reports}, 452(6), 155--176.

\bibitem[Peskin \& Schroeder(1995)]{HiggsVEV}
Peskin, M.~E. \& Schroeder, D.~V. (1995).
\newblock \textit{An Introduction to Quantum Field Theory}.
\newblock Addison-Wesley. (Ch. 20, electroweak symmetry breaking).

\bibitem[Polettini(2013)]{Polettini2013}
Polettini, M. (2013).
\newblock Fact-checking Ziegler's maximum entropy production principle beyond the linear regime and towards steady states.
\newblock \textit{Entropy}, 15(7), 2570--2584.
\newblock DOI: 10.3390/e15072570.

\bibitem[Verlinde(2011)]{Verlinde2011}
Verlinde, E. (2011).
\newblock On the origin of gravity and the laws of Newton.
\newblock \textit{JHEP}, 04, 029.
\newblock DOI: 10.1007/JHEP04(2011)029.

\bibitem[Weinberg(1987)]{Weinberg1987}
Weinberg, S. (1987).
\newblock Anthropic bound on the cosmological constant.
\newblock \textit{Phys. Rev. Lett.}, 59, 2607.

\bibitem[Egan \& Lineweaver(2010)]{EntropyMaximum}
Egan, C.~A. \& Lineweaver, C.~H. (2010).
\newblock A larger estimate of the entropy of the universe.
\newblock \textit{ApJ}, 710, 1825.

\bibitem[Zurek(2003)]{Zurek2003}
Zurek, W.~H. (2003).
\newblock Decoherence, einselection, and the quantum origins of the classical.
\newblock \textit{Rev. Mod. Phys.}, 75, 715--775.

\bibitem[Zurek(2005)]{Zurek2005}
Zurek, W.~H. (2005).
\newblock Probabilities from entanglement, Born's rule from envariance.
\newblock \textit{Physical Review A}, 71, 052105. arXiv:quant-ph/0405161.
\newblock DOI: 10.1103/PhysRevA.71.052105.

\bibitem[Schlosshauer \& Fine(2005)]{SchlossFine2005}
Schlosshauer, M. \& Fine, A. (2005).
\newblock On Zurek's derivation of the Born rule.
\newblock \textit{Foundations of Physics}, 35, 197--213. arXiv:quant-ph/0312058.

\bibitem[Geiger(2026b)]{FSTII}
Geiger, L. (2026).
\newblock Functional Stability Theory II: Chemical Stability and Autocatalytic Selection.
\newblock \textit{Zenodo preprint}, 2026.
\newblock DOI: 10.5281/zenodo.20130563.

\bibitem[Geiger(2026c)]{GeigerRH2026c}
Geiger, L. (2026).
\newblock The Riemann Hypothesis Part III: The Analytical Rank-Finite Certificate.
\newblock \textit{Zenodo preprint}, March 2026.
\newblock Series Concept DOI: 10.5281/zenodo.19035640.

\bibitem[Geiger(2026e)]{GeigerRH2026b}
Geiger, L. (2026).
\newblock Even Dominance of the Weil Quadratic Form: Spectral Analysis, Computer-Assisted Proof, and the Resolvent Gap Theorem.
\newblock \textit{Zenodo preprint}, March 2026.
\newblock Concept DOI: 10.5281/zenodo.19035640.

\bibitem[Geiger(2026d)]{FST-Unified2026}
Geiger, L. (2026).
\newblock FST Spectrum Duality: The Renormalized Free Energy Principle as Physical Instantiation of Functional Stability Theory.
\newblock \textit{Zenodo preprint}, 2026.
\newblock DOI: 10.5281/zenodo.19036190.

\bibitem[Geiger(2026f)]{GeigerCRM2026}
Geiger, L. (2026).
\newblock The Curvature Relaxation Model: A Four-Paper Program for Geometric Cosmology Without the Dark Sector.
\newblock \textit{Zenodo preprint}, 2026.
\newblock Concept DOI: 10.5281/zenodo.18728935.

\bibitem[Nash(1950)]{Nash1950}
Nash, J.~F. (1950).
\newblock Equilibrium points in $n$-person games.
\newblock \textit{Proceedings of the National Academy of Sciences}, 36(1), 48--49.

\bibitem[Paltridge(1975)]{Paltridge1975}
Paltridge, G.~W. (1975).
\newblock Global dynamics and climate---a system of minimum entropy exchange.
\newblock \textit{Quarterly Journal of the Royal Meteorological Society}, 101(429), 475--484.

\bibitem[Prigogine(1967)]{Prigogine1967}
Prigogine, I. (1967).
\newblock \textit{Introduction to Thermodynamics of Irreversible Processes} (3rd ed.).
\newblock Wiley Interscience.

\bibitem[Sawada et al.(2025)]{Sawada2025}
Sawada, Y., Daigaku, Y. \& Toma, K. (2025).
\newblock Maximum entropy production principle of thermodynamics for the birth and evolution of life.
\newblock \textit{Entropy}, 27(4), 449.
\newblock DOI: 10.3390/e27040449.

\end{thebibliography}

\medskip
\noindent\textbf{KI-Offenlegung.} Diese Arbeit wurde von Claude (Anthropic) bei der Literatursynthese, mathematischen Verifikation und Texterstellung unterst\"utzt. Alle wissenschaftlichen Behauptungen, originalen Argumente und theoretischen Beitr\"age liegen in der alleinigen Verantwortung des menschlichen Autors. Kein KI-System ist als Koautor gelistet.
}

\end{document}

